% Options for packages loaded elsewhere
\PassOptionsToPackage{unicode}{hyperref}
\PassOptionsToPackage{hyphens}{url}
\documentclass[
  12pt,
]{ctexart}
\usepackage{xcolor}
\usepackage{amsmath,amssymb}
\setcounter{secnumdepth}{-\maxdimen} % remove section numbering
\usepackage{iftex}
\ifPDFTeX
  \usepackage[T1]{fontenc}
  \usepackage[utf8]{inputenc}
  \usepackage{textcomp} % provide euro and other symbols
\else % if luatex or xetex
  \usepackage{unicode-math} % this also loads fontspec
  \defaultfontfeatures{Scale=MatchLowercase}
  \defaultfontfeatures[\rmfamily]{Ligatures=TeX,Scale=1}
\fi
\usepackage{lmodern}
\ifPDFTeX\else
  % xetex/luatex font selection
\fi
% Use upquote if available, for straight quotes in verbatim environments
\IfFileExists{upquote.sty}{\usepackage{upquote}}{}
\IfFileExists{microtype.sty}{% use microtype if available
  \usepackage[]{microtype}
  \UseMicrotypeSet[protrusion]{basicmath} % disable protrusion for tt fonts
}{}
\makeatletter
\@ifundefined{KOMAClassName}{% if non-KOMA class
  \IfFileExists{parskip.sty}{%
    \usepackage{parskip}
  }{% else
    \setlength{\parindent}{0pt}
    \setlength{\parskip}{6pt plus 2pt minus 1pt}}
}{% if KOMA class
  \KOMAoptions{parskip=half}}
\makeatother
\usepackage{color}
\usepackage{fancyvrb}
\newcommand{\VerbBar}{|}
\newcommand{\VERB}{\Verb[commandchars=\\\{\}]}
\DefineVerbatimEnvironment{Highlighting}{Verbatim}{commandchars=\\\{\}}
% Add ',fontsize=\small' for more characters per line
\newenvironment{Shaded}{}{}
\newcommand{\AlertTok}[1]{\textcolor[rgb]{1.00,0.00,0.00}{\textbf{#1}}}
\newcommand{\AnnotationTok}[1]{\textcolor[rgb]{0.38,0.63,0.69}{\textbf{\textit{#1}}}}
\newcommand{\AttributeTok}[1]{\textcolor[rgb]{0.49,0.56,0.16}{#1}}
\newcommand{\BaseNTok}[1]{\textcolor[rgb]{0.25,0.63,0.44}{#1}}
\newcommand{\BuiltInTok}[1]{\textcolor[rgb]{0.00,0.50,0.00}{#1}}
\newcommand{\CharTok}[1]{\textcolor[rgb]{0.25,0.44,0.63}{#1}}
\newcommand{\CommentTok}[1]{\textcolor[rgb]{0.38,0.63,0.69}{\textit{#1}}}
\newcommand{\CommentVarTok}[1]{\textcolor[rgb]{0.38,0.63,0.69}{\textbf{\textit{#1}}}}
\newcommand{\ConstantTok}[1]{\textcolor[rgb]{0.53,0.00,0.00}{#1}}
\newcommand{\ControlFlowTok}[1]{\textcolor[rgb]{0.00,0.44,0.13}{\textbf{#1}}}
\newcommand{\DataTypeTok}[1]{\textcolor[rgb]{0.56,0.13,0.00}{#1}}
\newcommand{\DecValTok}[1]{\textcolor[rgb]{0.25,0.63,0.44}{#1}}
\newcommand{\DocumentationTok}[1]{\textcolor[rgb]{0.73,0.13,0.13}{\textit{#1}}}
\newcommand{\ErrorTok}[1]{\textcolor[rgb]{1.00,0.00,0.00}{\textbf{#1}}}
\newcommand{\ExtensionTok}[1]{#1}
\newcommand{\FloatTok}[1]{\textcolor[rgb]{0.25,0.63,0.44}{#1}}
\newcommand{\FunctionTok}[1]{\textcolor[rgb]{0.02,0.16,0.49}{#1}}
\newcommand{\ImportTok}[1]{\textcolor[rgb]{0.00,0.50,0.00}{\textbf{#1}}}
\newcommand{\InformationTok}[1]{\textcolor[rgb]{0.38,0.63,0.69}{\textbf{\textit{#1}}}}
\newcommand{\KeywordTok}[1]{\textcolor[rgb]{0.00,0.44,0.13}{\textbf{#1}}}
\newcommand{\NormalTok}[1]{#1}
\newcommand{\OperatorTok}[1]{\textcolor[rgb]{0.40,0.40,0.40}{#1}}
\newcommand{\OtherTok}[1]{\textcolor[rgb]{0.00,0.44,0.13}{#1}}
\newcommand{\PreprocessorTok}[1]{\textcolor[rgb]{0.74,0.48,0.00}{#1}}
\newcommand{\RegionMarkerTok}[1]{#1}
\newcommand{\SpecialCharTok}[1]{\textcolor[rgb]{0.25,0.44,0.63}{#1}}
\newcommand{\SpecialStringTok}[1]{\textcolor[rgb]{0.73,0.40,0.53}{#1}}
\newcommand{\StringTok}[1]{\textcolor[rgb]{0.25,0.44,0.63}{#1}}
\newcommand{\VariableTok}[1]{\textcolor[rgb]{0.10,0.09,0.49}{#1}}
\newcommand{\VerbatimStringTok}[1]{\textcolor[rgb]{0.25,0.44,0.63}{#1}}
\newcommand{\WarningTok}[1]{\textcolor[rgb]{0.38,0.63,0.69}{\textbf{\textit{#1}}}}
\usepackage{longtable,booktabs,array}
\usepackage{calc} % for calculating minipage widths
% Correct order of tables after \paragraph or \subparagraph
\usepackage{etoolbox}
\makeatletter
\patchcmd\longtable{\par}{\if@noskipsec\mbox{}\fi\par}{}{}
\makeatother
% Allow footnotes in longtable head/foot
\IfFileExists{footnotehyper.sty}{\usepackage{footnotehyper}}{\usepackage{footnote}}
\makesavenoteenv{longtable}
\setlength{\emergencystretch}{3em} % prevent overfull lines
\providecommand{\tightlist}{%
  \setlength{\itemsep}{0pt}\setlength{\parskip}{0pt}}
\usepackage{bookmark}
\IfFileExists{xurl.sty}{\usepackage{xurl}}{} % add URL line breaks if available
\urlstyle{same}
\hypersetup{
  hidelinks,
  pdfcreator={LaTeX via pandoc}}

\author{SCX}
\date{2026}

\begin{document}

\section{SCX
通用模块化架构设计}\label{scx-ux901aux7528ux6a21ux5757ux5316ux67b6ux6784ux8bbeux8ba1}

\begin{quote}
日期：2026-06-26 来源：基于 GPT
对话分析（\texttt{与gpt的对话202606261332.md}）
目标：解决''每加一个领域就要写新的 adapter，代码重复，接口不一致''的问题
\end{quote}

\begin{center}\rule{0.5\linewidth}{0.5pt}\end{center}

\subsection{目录}\label{ux76eeux5f55}

\begin{enumerate}
\def\labelenumi{\arabic{enumi}.}
\tightlist
\item
  \hyperref[1-ux67b6ux6784ux603bux89c8]{架构总览（ASCII 图）}
\item
  \hyperref[2-ux8bbeux8ba1ux539fux5219]{设计原则}
\item
  \hyperref[3-ux6838ux5fc3ux62bdux8c61ux63a5ux53e3ux5b9aux4e49]{核心抽象接口定义（ABC）}
\item
  \hyperref[4-scx-core-ux4e0dux53d8ux5c42]{SCX Core（不变层）}
\item
  \hyperref[5-scx-encoder-ux9886ux57dfux76f8ux5173ux4f46ux63a5ux53e3ux7edfux4e00]{SCX
  Encoder（领域相关，但接口统一）}
\item
  \hyperref[6-scx-adapter-ux81eaux52a8ux751fux6210]{SCX
  Adapter（自动生成）}
\item
  \hyperref[7-ux58f0ux660eux5f0fux9886ux57dfux914dux7f6e]{声明式领域配置（YAML/JSON）}
\item
  \hyperref[8-ux6587ux4ef6ux5939ux7ed3ux6784]{文件夹结构}
\item
  \hyperref[9-scx-domain-ux6ce8ux518cux4e0eux751fux547dux5468ux671f]{SCX
  Domain 注册与生命周期}
\item
  \hyperref[10-ux63d2ux4ef6ux7cfbux7edf]{插件系统（可选扩展）}
\item
  \hyperref[11-ux4e0eux73b0ux6709ux4ee3ux7801ux7684ux8fc1ux79fbux8defux5f84]{与现有代码的迁移路径}
\item
  \hyperref[12-ux65b0ux589eux9886ux57dfux7684ux6b65ux9aa4]{新增领域的步骤（3
  步）}
\item
  \hyperref[13-ux6570ux636eux6d41ux5168ux666f]{数据流全景}
\item
  \hyperref[14-ux9644ux5f55]{附录：领域配置参考}
\end{enumerate}

\begin{center}\rule{0.5\linewidth}{0.5pt}\end{center}

\subsection{1. 架构总览}\label{ux67b6ux6784ux603bux89c8}

\begin{verbatim}
┌─────────────────────────────────────────────────────────────────────────┐
│                        SCX Universal Core                               │
│                                                                         │
│  ┌─────────────┐  ┌──────────────┐  ┌─────────────┐  ┌─────────────┐   │
│  │ StateSpace  │  │ ExpertRelia- │  │ DataClassi- │  │ ActionPolicy│   │
│  │ S = {s1...K}│  │ bility       │  │ fier         │  │ a*(s)       │   │
│  │ R_m(s)      │  │ R_m(s), SCX  │  │ 4-classify   │  │ keep/route/ │   │
│  │ V(s)        │  │ → expert gap │  │ V/R/N/E     │  │ compress/…  │   │
│  └──────┬──────┘  └──────┬───────┘  └──────┬──────┘  └──────┬──────┘   │
│         │                │                 │                │          │
│  ┌──────┴────────────────┴─────────────────┴────────────────┴──────┐   │
│  │                    SCX State-Outcome Database                      │   │
│  │              (s, e, a, o) — persistent across domains              │   │
│  └────────────────────────────────────────────────────────────────────┘   │
└─────────────────────────────────────────────────────────────────────────┘
                                    │
                    ┌───────────────┴───────────────┐
                    │                               │
         ┌──────────▼──────────┐        ┌───────────▼───────────┐
         │    SCX Encoder      │        │   SCX Domain Config   │
         │   (plug per domain) │        │   (YAML declaration)  │
         │                     │        │                       │
         │ encode(x) → z      │        │ encoder: path        │
         │ distance(a,b) → d  │        │ experts: [...]       │
         │ cluster(Z) → s_k   │        │ state_config: {...}  │
         └──────────┬──────────┘        │ thresholds: {...}   │
                    │                   └───────────────────────┘
         ┌──────────▼──────────┐
         │  Domain Adapter     │
         │  (auto-generated)   │
         │                     │
         │ Encoder → StateSpace│
         │ → ExpertReliability │
         │ → DataClassifier    │
         └─────────────────────┘

Domain-specific (per industry):
┌──────────┐ ┌──────────┐ ┌──────────┐ ┌──────────┐ ┌──────────┐
│  MLIP    │ │  Health  │ │  CIFAR   │ │  Robot   │ │   FEM    │
│  ACE/SOAP│ │  ResNet  │ │  CNN     │ │Trajectory│ │  Mesh    │
└──────────┘ └──────────┘ └──────────┘ └──────────┘ └──────────┘

Plugins (optional):
┌──────────┐ ┌──────────┐ ┌──────────┐
│ SCX-     │ │ SCX-    │ │ SCX-     │
│ Compress │ │ Online   │ │Influence │
└──────────┘ └──────────┘ └──────────┘
\end{verbatim}

\subsubsection{层间依赖规则}\label{ux5c42ux95f4ux4f9dux8d56ux89c4ux5219}

\begin{verbatim}
SCX Core          ← 不依赖任何 domain 代码
SCX Encoder       ← 依赖 Core（实现抽象接口），不依赖其他 domain
SCX Domain Config ← 声明式，不依赖 Python 代码
SCX Adapter       ← 自动生成，桥接 Encoder 和 Core
SCX Plugins       ← 可选，可依赖 Core + Encoder + Adapter
\end{verbatim}

\textbf{核心原则}：Core 永远不知道任何 Encoder 或 Domain
的存在。所有领域知识通过 Encoder 接口 + Domain 配置注入。

\begin{center}\rule{0.5\linewidth}{0.5pt}\end{center}

\subsection{2. 设计原则}\label{ux8bbeux8ba1ux539fux5219}

\subsubsection{2.1 通用数学层}\label{ux901aux7528ux6570ux5b66ux5c42}

SCX 的数学核心不依赖任何领域概念：

\begin{longtable}[]{@{}
  >{\raggedright\arraybackslash}p{(\linewidth - 4\tabcolsep) * \real{0.2667}}
  >{\raggedright\arraybackslash}p{(\linewidth - 4\tabcolsep) * \real{0.2000}}
  >{\raggedright\arraybackslash}p{(\linewidth - 4\tabcolsep) * \real{0.5333}}@{}}
\toprule\noalign{}
\begin{minipage}[b]{\linewidth}\raggedright
数学量
\end{minipage} & \begin{minipage}[b]{\linewidth}\raggedright
符号
\end{minipage} & \begin{minipage}[b]{\linewidth}\raggedright
与领域无关的定义
\end{minipage} \\
\midrule\noalign{}
\endhead
\bottomrule\noalign{}
\endlastfoot
状态空间 & \texttt{S\ =\ \{s₁,\ ...,\ s\_K\}} & 表示空间的聚类划分 \\
专家期望风险 &
\texttt{R\_m(s)\ =\ E{[}ℓ(f\_m(x),\ f*(x))\ \textbar{}\ x\ ∈\ s{]}} &
任意损失函数下的条件期望 \\
专家可靠性 & \texttt{SCX\_m(s)\ =\ P(ℓ\ \textless{}\ τ\ \textbar{}\ s)}
& 阈值条件下的可接受概率 \\
数据价值 &
\texttt{V(s)\ =\ ρ(s)\ ·\ {[}1-D(s){]}\ ·\ A\_\{e*\}(s)\ ·\ W\_\{weak\}(s)\ ·\ {[}1-N(s){]}}
& 五维乘积 \\
数据分类 & \texttt{\{Valuable,\ Redundant,\ Noisy,\ ExpertDependent\}} &
基于阈值的四分类 \\
动作策略 & \texttt{a*(s)\ =\ argmax\_a\ E{[}U(a,s,e,o){]}\ -\ λC(a)} &
成本约束的效用最大化 \\
\end{longtable}

\subsubsection{2.2 控制反转（IoC）}\label{ux63a7ux5236ux53cdux8f6cioc}

领域代码不直接调用 Core。Core 通过抽象接口调用领域代码（Encoder）。

\begin{verbatim}
传统方式:  Domain → Core (领域代码主动调用框架)
SCX 方式:  Core → Encoder (框架通过抽象接口回调领域代码)
\end{verbatim}

\subsubsection{2.3
配置驱动而非代码驱动}\label{ux914dux7f6eux9a71ux52a8ux800cux975eux4ee3ux7801ux9a71ux52a8}

添加新领域不需要写 Python 集成代码，只需要： 1. 写一个 Encoder
类（实现接口） 2. 写一个 YAML 配置文件 3. 注册到 domain registry

\begin{center}\rule{0.5\linewidth}{0.5pt}\end{center}

\subsection{3.
核心抽象接口定义}\label{ux6838ux5fc3ux62bdux8c61ux63a5ux53e3ux5b9aux4e49}

\subsubsection{3.1 SCXStateEncoder}\label{scxstateencoder}

\begin{Shaded}
\begin{Highlighting}[]
\CommentTok{\# scx/encoders/base.py}

\ImportTok{from}\NormalTok{ abc }\ImportTok{import}\NormalTok{ ABC, abstractmethod}
\ImportTok{from}\NormalTok{ typing }\ImportTok{import}\NormalTok{ Any, Optional}
\ImportTok{import}\NormalTok{ numpy }\ImportTok{as}\NormalTok{ np}


\KeywordTok{class}\NormalTok{ SCXStateEncoder(ABC):}
    \CommentTok{"""将任意领域输入 x 编码为状态空间中的向量。}

\CommentTok{    所有领域（MLIP, Health, CIFAR, Robot, FEM, Semi, Text...）}
\CommentTok{    只需实现这三个方法，即可接入 SCX 完整框架。}
\CommentTok{    """}

    \AttributeTok{@abstractmethod}
    \KeywordTok{def}\NormalTok{ encode(}\VariableTok{self}\NormalTok{, x: Any) }\OperatorTok{{-}\textgreater{}}\NormalTok{ np.ndarray:}
        \CommentTok{"""将原始输入 x 映射为稠密向量 z ∈ R\^{}d。}

\CommentTok{        这是状态编码的核心。可以是：}
\CommentTok{        {-} MLIP: ACE/SOAP descriptor}
\CommentTok{        {-} Vision: CNN/ViT feature embedding}
\CommentTok{        {-} Trajectory: temporal encoding of (o, a, r) sequence}
\CommentTok{        {-} Tabular: normalized + embedded feature vector}
\CommentTok{        {-} Text: LLM embedding}
\CommentTok{        {-} Simulation: mesh/grid encoding}

\CommentTok{        Parameters}
\CommentTok{        {-}{-}{-}{-}{-}{-}{-}{-}{-}{-}}
\CommentTok{        x : Any}
\CommentTok{            单个原始输入样本。}

\CommentTok{        Returns}
\CommentTok{        {-}{-}{-}{-}{-}{-}{-}}
\CommentTok{        np.ndarray, shape (d,)}
\CommentTok{            编码后的稠密向量。}
\CommentTok{        """}
\NormalTok{        ...}

    \KeywordTok{def}\NormalTok{ encode\_batch(}\VariableTok{self}\NormalTok{, X: }\BuiltInTok{list}\NormalTok{[Any]) }\OperatorTok{{-}\textgreater{}}\NormalTok{ np.ndarray:}
        \CommentTok{"""批量编码，默认逐个调用 encode()，子类可重写做并行加速。"""}
        \ControlFlowTok{return}\NormalTok{ np.array([}\VariableTok{self}\NormalTok{.encode(x) }\ControlFlowTok{for}\NormalTok{ x }\KeywordTok{in}\NormalTok{ X])}

    \AttributeTok{@abstractmethod}
    \KeywordTok{def}\NormalTok{ distance(}\VariableTok{self}\NormalTok{, a: np.ndarray, b: np.ndarray) }\OperatorTok{{-}\textgreater{}} \BuiltInTok{float}\NormalTok{:}
        \CommentTok{"""定义状态空间中的距离度量。}

\CommentTok{        默认可以用欧氏距离，但不同领域可能需要自定义：}
\CommentTok{        {-} SOAP 特征用 cosine 距离}
\CommentTok{        {-} 轨迹用 DTW}
\CommentTok{        {-} 图结构用 graph edit distance}

\CommentTok{        Parameters}
\CommentTok{        {-}{-}{-}{-}{-}{-}{-}{-}{-}{-}}
\CommentTok{        a : np.ndarray, shape (d,)}
\CommentTok{        b : np.ndarray, shape (d,)}

\CommentTok{        Returns}
\CommentTok{        {-}{-}{-}{-}{-}{-}{-}}
\CommentTok{        float}
\CommentTok{            距离（越小越近）。}
\CommentTok{        """}
\NormalTok{        ...}

    \AttributeTok{@abstractmethod}
    \KeywordTok{def}\NormalTok{ cluster(}\VariableTok{self}\NormalTok{, Z: np.ndarray, }\OperatorTok{**}\NormalTok{kwargs) }\OperatorTok{{-}\textgreater{}} \BuiltInTok{tuple}\NormalTok{[np.ndarray, np.ndarray]:}
        \CommentTok{"""对编码后的向量集进行聚类，划分状态。}

\CommentTok{        Parameters}
\CommentTok{        {-}{-}{-}{-}{-}{-}{-}{-}{-}{-}}
\CommentTok{        Z : np.ndarray, shape (N, d)}
\CommentTok{            编码向量矩阵。}

\CommentTok{        Returns}
\CommentTok{        {-}{-}{-}{-}{-}{-}{-}}
\CommentTok{        labels : np.ndarray, shape (N,)}
\CommentTok{            每个样本的状态分配 \{0, ..., K{-}1\}。}
\CommentTok{        centroids : np.ndarray, shape (K, d)}
\CommentTok{            每个状态的质心。}
\CommentTok{        """}
\NormalTok{        ...}

    \KeywordTok{def}\NormalTok{ get\_feature\_dim(}\VariableTok{self}\NormalTok{) }\OperatorTok{{-}\textgreater{}} \BuiltInTok{int}\NormalTok{:}
        \CommentTok{"""返回编码向量的维度 d。"""}
        \ControlFlowTok{if} \BuiltInTok{hasattr}\NormalTok{(}\VariableTok{self}\NormalTok{, }\StringTok{\textquotesingle{}\_feature\_dim\textquotesingle{}}\NormalTok{):}
            \ControlFlowTok{return} \VariableTok{self}\NormalTok{.\_feature\_dim}
        \CommentTok{\# 尝试通过一个 dummy 输入推断}
\NormalTok{        dummy }\OperatorTok{=} \VariableTok{self}\NormalTok{.encode(}\VariableTok{self}\NormalTok{.\_dummy\_input()) }\ControlFlowTok{if} \BuiltInTok{hasattr}\NormalTok{(}\VariableTok{self}\NormalTok{, }\StringTok{\textquotesingle{}\_dummy\_input\textquotesingle{}}\NormalTok{) }\ControlFlowTok{else} \VariableTok{None}
        \ControlFlowTok{if}\NormalTok{ dummy }\KeywordTok{is} \KeywordTok{not} \VariableTok{None}\NormalTok{:}
            \ControlFlowTok{return}\NormalTok{ dummy.shape[}\DecValTok{0}\NormalTok{]}
        \ControlFlowTok{return} \DecValTok{0}

    \KeywordTok{def}\NormalTok{ \_dummy\_input(}\VariableTok{self}\NormalTok{) }\OperatorTok{{-}\textgreater{}}\NormalTok{ Any:}
        \CommentTok{"""子类可重写以提供用于推断维度的模拟输入。"""}
        \ControlFlowTok{raise} \PreprocessorTok{NotImplementedError}
\end{Highlighting}
\end{Shaded}

\subsubsection{3.2
SCXExpert（轻量抽象）}\label{scxexpertux8f7bux91cfux62bdux8c61}

\begin{Shaded}
\begin{Highlighting}[]
\CommentTok{\# scx/core/expert.py}

\ImportTok{from}\NormalTok{ abc }\ImportTok{import}\NormalTok{ ABC, abstractmethod}
\ImportTok{from}\NormalTok{ typing }\ImportTok{import}\NormalTok{ Any}


\KeywordTok{class}\NormalTok{ SCXExpert(ABC):}
    \CommentTok{"""任意专家：ML 模型、人、仿真器、API、规则系统。}

\CommentTok{    核心思想：SCX 不关心专家内部实现，只关心：}
\CommentTok{    {-} predict(x) → y：预测输出}
\CommentTok{    {-} cost() → float：使用成本（用于成本约束决策）}
\CommentTok{    """}

    \AttributeTok{@abstractmethod}
    \KeywordTok{def}\NormalTok{ predict(}\VariableTok{self}\NormalTok{, x: Any) }\OperatorTok{{-}\textgreater{}}\NormalTok{ Any:}
        \CommentTok{"""对单个输入 x 做出预测/决策。}

\CommentTok{        Parameters}
\CommentTok{        {-}{-}{-}{-}{-}{-}{-}{-}{-}{-}}
\CommentTok{        x : Any}
\CommentTok{            输入样本。}

\CommentTok{        Returns}
\CommentTok{        {-}{-}{-}{-}{-}{-}{-}}
\CommentTok{        Any}
\CommentTok{            专家输出（可以是标量、向量、类别、文本等）。}
\CommentTok{        """}
\NormalTok{        ...}

    \KeywordTok{def}\NormalTok{ predict\_batch(}\VariableTok{self}\NormalTok{, X: }\BuiltInTok{list}\NormalTok{[Any]) }\OperatorTok{{-}\textgreater{}} \BuiltInTok{list}\NormalTok{[Any]:}
        \CommentTok{"""批量预测，默认逐个调用。"""}
        \ControlFlowTok{return}\NormalTok{ [}\VariableTok{self}\NormalTok{.predict(x) }\ControlFlowTok{for}\NormalTok{ x }\KeywordTok{in}\NormalTok{ X]}

    \AttributeTok{@abstractmethod}
    \KeywordTok{def}\NormalTok{ cost(}\VariableTok{self}\NormalTok{) }\OperatorTok{{-}\textgreater{}} \BuiltInTok{float}\NormalTok{:}
        \CommentTok{"""返回使用该专家一次的成本（相对值，用于路由决策）。}

\CommentTok{        例如：}
\CommentTok{        {-} 轻量模型: 0.1}
\CommentTok{        {-} DFT 计算: 100.0}
\CommentTok{        {-} 人类专家: 500.0}
\CommentTok{        {-} 高保真仿真: 50.0}
\CommentTok{        """}
\NormalTok{        ...}

    \AttributeTok{@property}
    \KeywordTok{def}\NormalTok{ name(}\VariableTok{self}\NormalTok{) }\OperatorTok{{-}\textgreater{}} \BuiltInTok{str}\NormalTok{:}
        \CommentTok{"""专家名称，用于日志和报告。"""}
        \ControlFlowTok{return} \VariableTok{self}\NormalTok{.}\VariableTok{\_\_class\_\_}\NormalTok{.}\VariableTok{\_\_name\_\_}

    \AttributeTok{@property}
    \KeywordTok{def}\NormalTok{ metadata(}\VariableTok{self}\NormalTok{) }\OperatorTok{{-}\textgreater{}} \BuiltInTok{dict}\NormalTok{:}
        \CommentTok{"""专家元数据（可选的扩展信息）。"""}
        \ControlFlowTok{return}\NormalTok{ \{\}}
\end{Highlighting}
\end{Shaded}

\subsubsection{3.3
SCXDataPoint（可选抽象）}\label{scxdatapointux53efux9009ux62bdux8c61}

\begin{Shaded}
\begin{Highlighting}[]
\CommentTok{\# scx/core/datapoint.py}

\ImportTok{from}\NormalTok{ abc }\ImportTok{import}\NormalTok{ ABC, abstractmethod}
\ImportTok{from}\NormalTok{ typing }\ImportTok{import}\NormalTok{ Any}
\ImportTok{import}\NormalTok{ numpy }\ImportTok{as}\NormalTok{ np}


\KeywordTok{class}\NormalTok{ SCXDataPoint(ABC):}
    \CommentTok{"""任意数据点：DFT 帧、图像、轨迹片段、仿真结果、文本。}

\CommentTok{    方便封装常见数据操作。非必须——可以直接传原始类型给 Encoder。}
\CommentTok{    """}

    \AttributeTok{@abstractmethod}
    \KeywordTok{def}\NormalTok{ get\_features(}\VariableTok{self}\NormalTok{) }\OperatorTok{{-}\textgreater{}}\NormalTok{ np.ndarray:}
        \CommentTok{"""返回该数据点的特征向量（用于编码前处理）。"""}
\NormalTok{        ...}

    \KeywordTok{def}\NormalTok{ get\_label(}\VariableTok{self}\NormalTok{) }\OperatorTok{{-}\textgreater{}}\NormalTok{ Any:}
        \CommentTok{"""返回该数据点的标签（如果可用）。"""}
        \ControlFlowTok{return} \VariableTok{None}

    \AttributeTok{@classmethod}
    \AttributeTok{@abstractmethod}
    \KeywordTok{def}\NormalTok{ from\_raw(cls, raw: Any) }\OperatorTok{{-}\textgreater{}} \StringTok{"SCXDataPoint"}\NormalTok{:}
        \CommentTok{"""从原始格式创建数据点。"""}
\NormalTok{        ...}
\end{Highlighting}
\end{Shaded}

\begin{center}\rule{0.5\linewidth}{0.5pt}\end{center}

\subsection{4. SCX Core（不变层）}\label{scx-coreux4e0dux53d8ux5c42}

Core 层完全保持通用，不引用任何领域代码。它包括：

\subsubsection{4.1
现有模块（保留，不重写）}\label{ux73b0ux6709ux6a21ux5757ux4fddux7559ux4e0dux91cdux5199}

\begin{longtable}[]{@{}
  >{\raggedright\arraybackslash}p{(\linewidth - 4\tabcolsep) * \real{0.3333}}
  >{\raggedright\arraybackslash}p{(\linewidth - 4\tabcolsep) * \real{0.3333}}
  >{\raggedright\arraybackslash}p{(\linewidth - 4\tabcolsep) * \real{0.3333}}@{}}
\toprule\noalign{}
\begin{minipage}[b]{\linewidth}\raggedright
模块
\end{minipage} & \begin{minipage}[b]{\linewidth}\raggedright
文件
\end{minipage} & \begin{minipage}[b]{\linewidth}\raggedright
说明
\end{minipage} \\
\midrule\noalign{}
\endhead
\bottomrule\noalign{}
\endlastfoot
\texttt{SCXFramework} & \texttt{core/framework.py} &
顶层编排类，重构为接收 Encoder + Config \\
\texttt{SCXConfig} & \texttt{core/config.py} & 扩展为支持 domain config
加载 \\
\texttt{StateSpace} & \texttt{state/space.py} & 状态空间容器（不变） \\
\texttt{StateDiscovery} & \texttt{state/discovery.py} & 状态发现（接收
Encoder.cluster） \\
\texttt{ExpertReliability} & \texttt{expert/reliability.py} &
专家可靠性估计（不变） \\
\texttt{ExpertRegistry} & \texttt{expert/registry.py} &
专家注册（不变） \\
\texttt{DataClassifier} & \texttt{valuation/classifier.py} &
四分类（不变） \\
\texttt{ActionPolicy} & \texttt{action/policy.py} & 动作策略（不变） \\
\texttt{OnlineSCXFramework} & \texttt{core/online.py} &
在线监测（不变） \\
\end{longtable}

\subsubsection{4.2 新增 Core 模块}\label{ux65b0ux589e-core-ux6a21ux5757}

\paragraph{SCXDomainEngine}\label{scxdomainengine}

\begin{Shaded}
\begin{Highlighting}[]
\CommentTok{\# scx/core/domain\_engine.py}

\KeywordTok{class}\NormalTok{ SCXDomainEngine:}
    \CommentTok{"""领域引擎：将 Encoder + Config + Core 串联的胶水层。}

\CommentTok{    这个类是用户的主要入口点。替代原来直接使用 SCXFramework。}
\CommentTok{    """}

    \KeywordTok{def} \FunctionTok{\_\_init\_\_}\NormalTok{(}\VariableTok{self}\NormalTok{, domain\_config: }\BuiltInTok{dict} \OperatorTok{|} \BuiltInTok{str}\NormalTok{):}
        \CommentTok{\# 加载 YAML/JSON 配置}
        \CommentTok{\# 动态实例化 Encoder}
        \CommentTok{\# 创建 ExpertRegistry}
        \CommentTok{\# 配置阈值}
\NormalTok{        ...}

    \KeywordTok{def}\NormalTok{ fit(}\VariableTok{self}\NormalTok{, X, y}\OperatorTok{=}\VariableTok{None}\NormalTok{):}
        \CommentTok{"""运行完整 SCX pipeline。"""}

    \KeywordTok{def}\NormalTok{ audit(}\VariableTok{self}\NormalTok{) }\OperatorTok{{-}\textgreater{}} \StringTok{"SCXAuditReport"}\NormalTok{:}
        \CommentTok{"""生成数据价值审计报告。"""}

    \KeywordTok{def}\NormalTok{ value(}\VariableTok{self}\NormalTok{, x) }\OperatorTok{{-}\textgreater{}} \BuiltInTok{float}\NormalTok{:}
        \CommentTok{"""计算单个样本的数据价值。"""}

    \KeywordTok{def}\NormalTok{ recommend\_action(}\VariableTok{self}\NormalTok{, state\_id) }\OperatorTok{{-}\textgreater{}} \BuiltInTok{str}\NormalTok{:}
        \CommentTok{"""推荐动作。"""}

    \KeywordTok{def}\NormalTok{ save(}\VariableTok{self}\NormalTok{, path): ...}
    \KeywordTok{def}\NormalTok{ load(}\VariableTok{self}\NormalTok{, path): ...}
\end{Highlighting}
\end{Shaded}

\paragraph{SCXAuditReport}\label{scxauditreport}

\begin{Shaded}
\begin{Highlighting}[]
\CommentTok{\# scx/core/audit.py}

\KeywordTok{class}\NormalTok{ SCXAuditReport:}
    \CommentTok{"""标准审计报告，跨领域统一格式。}

\CommentTok{    报告包含：}
\CommentTok{    1. 状态覆盖地图}
\CommentTok{    2. 冗余比例}
\CommentTok{    3. 噪声风险}
\CommentTok{    4. 高价值状态}
\CommentTok{    5. Expert 可靠性矩阵}
\CommentTok{    6. 建议动作}
\CommentTok{    7. 执行后收益估计}
\CommentTok{    """}
\end{Highlighting}
\end{Shaded}

\begin{center}\rule{0.5\linewidth}{0.5pt}\end{center}

\subsection{5. SCX
Encoder（领域相关，但接口统一）}\label{scx-encoderux9886ux57dfux76f8ux5173ux4f46ux63a5ux53e3ux7edfux4e00}

\subsubsection{5.1 内置 Encoder
列表}\label{ux5185ux7f6e-encoder-ux5217ux8868}

\begin{verbatim}
scx/encoders/
├── base.py              # SCXStateEncoder 抽象基类
├── mlip.py              # ACE/SOAP descriptor → 状态向量
├── vision.py            # CNN/ViT/CLIP → 图像 embedding
├── tabular.py           # 表格数据 → 归一化向量
├── trajectory.py        # 机器人轨迹 → temporal encoding
├── simulation.py        # FEM/PDE mesh → response encoding
├── text.py              # LLM embedding → 语义向量
└── health.py            # 医学图像 → ResNet feature
\end{verbatim}

\subsubsection{5.2 Encoder 示例}\label{encoder-ux793aux4f8b}

\paragraph{MLIP Encoder}\label{mlip-encoder}

\begin{Shaded}
\begin{Highlighting}[]
\CommentTok{\# scx/encoders/mlip.py}

\KeywordTok{class}\NormalTok{ ACEEncoder(SCXStateEncoder):}
    \CommentTok{"""ACE/SOAP descriptor 编码器，用于原子构型。"""}

    \KeywordTok{def} \FunctionTok{\_\_init\_\_}\NormalTok{(}\VariableTok{self}\NormalTok{, descriptor: }\BuiltInTok{str} \OperatorTok{=} \StringTok{"soap"}\NormalTok{, rcut: }\BuiltInTok{float} \OperatorTok{=} \FloatTok{5.0}\NormalTok{, nmax: }\BuiltInTok{int} \OperatorTok{=} \DecValTok{8}\NormalTok{, lmax: }\BuiltInTok{int} \OperatorTok{=} \DecValTok{6}\NormalTok{):}
        \VariableTok{self}\NormalTok{.descriptor\_type }\OperatorTok{=}\NormalTok{ descriptor}
        \VariableTok{self}\NormalTok{.rcut }\OperatorTok{=}\NormalTok{ rcut}
        \VariableTok{self}\NormalTok{.nmax }\OperatorTok{=}\NormalTok{ nmax}
        \VariableTok{self}\NormalTok{.lmax }\OperatorTok{=}\NormalTok{ lmax}
        \CommentTok{\# lazily initialize dscribe/etc. here}

    \KeywordTok{def}\NormalTok{ encode(}\VariableTok{self}\NormalTok{, x: }\StringTok{"ase.Atoms"}\NormalTok{) }\OperatorTok{{-}\textgreater{}}\NormalTok{ np.ndarray:}
        \CommentTok{\# x 是 ASE Atoms 对象}
        \CommentTok{\# 返回 SOAP/ACE 描述符向量}
\NormalTok{        ...}

    \KeywordTok{def}\NormalTok{ distance(}\VariableTok{self}\NormalTok{, a: np.ndarray, b: np.ndarray) }\OperatorTok{{-}\textgreater{}} \BuiltInTok{float}\NormalTok{:}
        \CommentTok{\# SOAP 常用 cosine 距离}
        \ControlFlowTok{return} \FloatTok{1.0} \OperatorTok{{-}}\NormalTok{ np.dot(a, b) }\OperatorTok{/}\NormalTok{ (np.linalg.norm(a) }\OperatorTok{*}\NormalTok{ np.linalg.norm(b) }\OperatorTok{+} \FloatTok{1e{-}10}\NormalTok{)}

    \KeywordTok{def}\NormalTok{ cluster(}\VariableTok{self}\NormalTok{, Z: np.ndarray, }\OperatorTok{**}\NormalTok{kwargs) }\OperatorTok{{-}\textgreater{}} \BuiltInTok{tuple}\NormalTok{[np.ndarray, np.ndarray]:}
        \CommentTok{\# 用 HDBSCAN 或 K{-}Means 聚类}
        \ImportTok{from}\NormalTok{ sklearn.cluster }\ImportTok{import}\NormalTok{ KMeans}
\NormalTok{        n }\OperatorTok{=}\NormalTok{ kwargs.get(}\StringTok{"n\_states"}\NormalTok{, }\DecValTok{10}\NormalTok{)}
\NormalTok{        km }\OperatorTok{=}\NormalTok{ KMeans(n\_clusters}\OperatorTok{=}\NormalTok{n, random\_state}\OperatorTok{=}\DecValTok{42}\NormalTok{)}
\NormalTok{        labels }\OperatorTok{=}\NormalTok{ km.fit\_predict(Z)}
        \ControlFlowTok{return}\NormalTok{ labels, km.cluster\_centers\_}
\end{Highlighting}
\end{Shaded}

\paragraph{Vision Encoder}\label{vision-encoder}

\begin{Shaded}
\begin{Highlighting}[]
\CommentTok{\# scx/encoders/vision.py}

\KeywordTok{class}\NormalTok{ CNNEncoder(SCXStateEncoder):}
    \CommentTok{"""CNN/ViT encoder for images."""}

    \KeywordTok{def} \FunctionTok{\_\_init\_\_}\NormalTok{(}\VariableTok{self}\NormalTok{, backbone: }\BuiltInTok{str} \OperatorTok{=} \StringTok{"resnet18"}\NormalTok{, device: }\BuiltInTok{str} \OperatorTok{=} \StringTok{"cpu"}\NormalTok{):}
        \ImportTok{import}\NormalTok{ torch}
        \VariableTok{self}\NormalTok{.device }\OperatorTok{=}\NormalTok{ device}
        \VariableTok{self}\NormalTok{.model }\OperatorTok{=} \VariableTok{self}\NormalTok{.\_build\_backbone(backbone).to(device)}
        \VariableTok{self}\NormalTok{.model.}\BuiltInTok{eval}\NormalTok{()}

    \KeywordTok{def}\NormalTok{ encode(}\VariableTok{self}\NormalTok{, x: }\StringTok{"torch.Tensor | np.ndarray"}\NormalTok{) }\OperatorTok{{-}\textgreater{}}\NormalTok{ np.ndarray:}
        \ImportTok{import}\NormalTok{ torch}
        \ControlFlowTok{if} \BuiltInTok{isinstance}\NormalTok{(x, np.ndarray):}
\NormalTok{            x }\OperatorTok{=}\NormalTok{ torch.from\_numpy(x).to(}\VariableTok{self}\NormalTok{.device)}
        \ControlFlowTok{if}\NormalTok{ x.ndim }\OperatorTok{==} \DecValTok{3}\NormalTok{:}
\NormalTok{            x }\OperatorTok{=}\NormalTok{ x.unsqueeze(}\DecValTok{0}\NormalTok{)  }\CommentTok{\# add batch dim}
        \ControlFlowTok{with}\NormalTok{ torch.no\_grad():}
\NormalTok{            feat }\OperatorTok{=} \VariableTok{self}\NormalTok{.model.forward\_features(x)}
        \ControlFlowTok{return}\NormalTok{ feat.cpu().numpy().flatten()}

    \KeywordTok{def}\NormalTok{ distance(}\VariableTok{self}\NormalTok{, a: np.ndarray, b: np.ndarray) }\OperatorTok{{-}\textgreater{}} \BuiltInTok{float}\NormalTok{:}
        \ControlFlowTok{return} \BuiltInTok{float}\NormalTok{(np.linalg.norm(a }\OperatorTok{{-}}\NormalTok{ b))}

    \KeywordTok{def}\NormalTok{ cluster(}\VariableTok{self}\NormalTok{, Z: np.ndarray, }\OperatorTok{**}\NormalTok{kwargs) }\OperatorTok{{-}\textgreater{}} \BuiltInTok{tuple}\NormalTok{[np.ndarray, np.ndarray]:}
        \ImportTok{from}\NormalTok{ sklearn.cluster }\ImportTok{import}\NormalTok{ MiniBatchKMeans}
\NormalTok{        n }\OperatorTok{=}\NormalTok{ kwargs.get(}\StringTok{"n\_states"}\NormalTok{, }\DecValTok{20}\NormalTok{)}
\NormalTok{        km }\OperatorTok{=}\NormalTok{ MiniBatchKMeans(n\_clusters}\OperatorTok{=}\NormalTok{n, random\_state}\OperatorTok{=}\DecValTok{42}\NormalTok{)}
\NormalTok{        labels }\OperatorTok{=}\NormalTok{ km.fit\_predict(Z)}
        \ControlFlowTok{return}\NormalTok{ labels, km.cluster\_centers\_}
\end{Highlighting}
\end{Shaded}

\paragraph{Trajectory Encoder}\label{trajectory-encoder}

\begin{Shaded}
\begin{Highlighting}[]
\CommentTok{\# scx/encoders/trajectory.py}

\KeywordTok{class}\NormalTok{ TrajectoryEncoder(SCXStateEncoder):}
    \CommentTok{"""机器人轨迹编码器。"""}

    \KeywordTok{def} \FunctionTok{\_\_init\_\_}\NormalTok{(}\VariableTok{self}\NormalTok{, state\_dim: }\BuiltInTok{int} \OperatorTok{=} \DecValTok{128}\NormalTok{, use\_temporal: }\BuiltInTok{bool} \OperatorTok{=} \VariableTok{True}\NormalTok{):}
        \VariableTok{self}\NormalTok{.state\_dim }\OperatorTok{=}\NormalTok{ state\_dim}
        \VariableTok{self}\NormalTok{.use\_temporal }\OperatorTok{=}\NormalTok{ use\_temporal}
        \CommentTok{\# lazily build encoder network}

    \KeywordTok{def}\NormalTok{ encode(}\VariableTok{self}\NormalTok{, x: }\BuiltInTok{dict}\NormalTok{) }\OperatorTok{{-}\textgreater{}}\NormalTok{ np.ndarray:}
        \CommentTok{\# x 包含 observation, action, reward, success 等}
        \CommentTok{\# 返回固定维度的轨迹 embedding}
\NormalTok{        ...}

    \KeywordTok{def}\NormalTok{ distance(}\VariableTok{self}\NormalTok{, a: np.ndarray, b: np.ndarray) }\OperatorTok{{-}\textgreater{}} \BuiltInTok{float}\NormalTok{:}
        \CommentTok{\# 轨迹用 DTW 或简单欧氏距离}
        \ControlFlowTok{return} \BuiltInTok{float}\NormalTok{(np.linalg.norm(a }\OperatorTok{{-}}\NormalTok{ b))}

    \KeywordTok{def}\NormalTok{ cluster(}\VariableTok{self}\NormalTok{, Z: np.ndarray, }\OperatorTok{**}\NormalTok{kwargs) }\OperatorTok{{-}\textgreater{}} \BuiltInTok{tuple}\NormalTok{[np.ndarray, np.ndarray]:}
        \ImportTok{from}\NormalTok{ sklearn.cluster }\ImportTok{import}\NormalTok{ HDBSCAN}
\NormalTok{        clusterer }\OperatorTok{=}\NormalTok{ HDBSCAN(min\_cluster\_size}\OperatorTok{=}\NormalTok{kwargs.get(}\StringTok{"min\_cluster\_size"}\NormalTok{, }\DecValTok{5}\NormalTok{))}
\NormalTok{        labels }\OperatorTok{=}\NormalTok{ clusterer.fit\_predict(Z)}
        \CommentTok{\# HDBSCAN 可能产生 {-}1 (noise)，需要处理}
\NormalTok{        centroids }\OperatorTok{=}\NormalTok{ ...}
        \ControlFlowTok{return}\NormalTok{ labels, centroids}
\end{Highlighting}
\end{Shaded}

\begin{center}\rule{0.5\linewidth}{0.5pt}\end{center}

\subsection{6. SCX
Adapter（自动生成）}\label{scx-adapterux81eaux52a8ux751fux6210}

Adapter 是 Encoder 和 Core 之间的桥接层。它从 Encoder 和 Domain Config
自动推导出 Core 所需的全部组件。

\subsubsection{6.1
自动适配的内容}\label{ux81eaux52a8ux9002ux914dux7684ux5185ux5bb9}

\begin{longtable}[]{@{}
  >{\raggedright\arraybackslash}p{(\linewidth - 4\tabcolsep) * \real{0.2400}}
  >{\raggedright\arraybackslash}p{(\linewidth - 4\tabcolsep) * \real{0.2400}}
  >{\raggedright\arraybackslash}p{(\linewidth - 4\tabcolsep) * \real{0.5200}}@{}}
\toprule\noalign{}
\begin{minipage}[b]{\linewidth}\raggedright
组件
\end{minipage} & \begin{minipage}[b]{\linewidth}\raggedright
来源
\end{minipage} & \begin{minipage}[b]{\linewidth}\raggedright
自动适配逻辑
\end{minipage} \\
\midrule\noalign{}
\endhead
\bottomrule\noalign{}
\endlastfoot
\texttt{StateSpace} & Encoder.cluster() 结果 & 从 labels + centroids
构建 \\
Expert 列表 & Domain Config \texttt{experts} 字段 & 根据 type
自动实例化专家 \\
损失函数 & 默认 MSE 或 Config 指定 & 根据问题类型选择 \\
\texttt{R\_m(s)} & ExpertReliability.estimate() & 使用 Encoder 的 state
assignment \\
阈值 & Domain Config \texttt{thresholds} & 覆盖 Core 默认值 \\
特征维度 & Encoder.get\_feature\_dim() & 自动推断 \\
聚类方法 & Encoder.cluster() & 直接委托 \\
\end{longtable}

\subsubsection{6.2 Adapter
生成伪代码}\label{adapter-ux751fux6210ux4f2aux4ee3ux7801}

\begin{Shaded}
\begin{Highlighting}[]
\CommentTok{\# scx/core/adapter.py}

\KeywordTok{def}\NormalTok{ build\_from\_domain(domain\_config: }\BuiltInTok{dict}\NormalTok{) }\OperatorTok{{-}\textgreater{}}\NormalTok{ SCXDomainEngine:}
    \CommentTok{"""从领域配置自动构建 SCXDomainEngine。"""}

    \CommentTok{\# 1. 实例化 Encoder}
\NormalTok{    encoder\_cls }\OperatorTok{=}\NormalTok{ import\_class(domain\_config[}\StringTok{"encoder"}\NormalTok{])}
\NormalTok{    encoder }\OperatorTok{=}\NormalTok{ encoder\_cls(}\OperatorTok{**}\NormalTok{domain\_config.get(}\StringTok{"encoder\_params"}\NormalTok{, \{\}))}

    \CommentTok{\# 2. 创建 ExpertRegistry}
\NormalTok{    registry }\OperatorTok{=}\NormalTok{ ExpertRegistry()}
    \ControlFlowTok{for}\NormalTok{ exp\_cfg }\KeywordTok{in}\NormalTok{ domain\_config.get(}\StringTok{"experts"}\NormalTok{, []):}
\NormalTok{        expert }\OperatorTok{=}\NormalTok{ create\_expert(exp\_cfg)  }\CommentTok{\# 根据 type 工厂方法}
\NormalTok{        registry.register(expert)}

    \CommentTok{\# 3. 构建 SCXConfig（叠加领域阈值）}
\NormalTok{    base\_config }\OperatorTok{=}\NormalTok{ SCXConfig()}
\NormalTok{    overrides }\OperatorTok{=}\NormalTok{ domain\_config.get(}\StringTok{"thresholds"}\NormalTok{, \{\})}
    \ControlFlowTok{for}\NormalTok{ k, v }\KeywordTok{in}\NormalTok{ overrides.items():}
        \ControlFlowTok{if} \BuiltInTok{hasattr}\NormalTok{(base\_config, k):}
            \BuiltInTok{setattr}\NormalTok{(base\_config, k, v)}
\NormalTok{    config }\OperatorTok{=}\NormalTok{ base\_config}

    \CommentTok{\# 4. 构建领域引擎}
\NormalTok{    engine }\OperatorTok{=}\NormalTok{ SCXDomainEngine(config}\OperatorTok{=}\NormalTok{config, encoder}\OperatorTok{=}\NormalTok{encoder, registry}\OperatorTok{=}\NormalTok{registry)}
    \ControlFlowTok{return}\NormalTok{ engine}
\end{Highlighting}
\end{Shaded}

\subsubsection{6.3 Encoder → StateSpace
的自动推导}\label{encoder-statespace-ux7684ux81eaux52a8ux63a8ux5bfc}

\begin{Shaded}
\begin{Highlighting}[]
\CommentTok{\# SCXDomainEngine.fit() 内部}

\KeywordTok{def}\NormalTok{ fit(}\VariableTok{self}\NormalTok{, X, y}\OperatorTok{=}\VariableTok{None}\NormalTok{):}
    \CommentTok{\# 1. Encode}
\NormalTok{    Z }\OperatorTok{=} \VariableTok{self}\NormalTok{.encoder.encode\_batch(X)}

    \CommentTok{\# 2. Cluster → states}
\NormalTok{    labels, centroids }\OperatorTok{=} \VariableTok{self}\NormalTok{.encoder.cluster(Z, n\_states}\OperatorTok{=}\VariableTok{self}\NormalTok{.config.n\_states)}

    \CommentTok{\# 3. Build StateSpace}
    \VariableTok{self}\NormalTok{.state\_space }\OperatorTok{=}\NormalTok{ StateSpace(n\_states}\OperatorTok{=}\BuiltInTok{len}\NormalTok{(centroids))}
    \VariableTok{self}\NormalTok{.state\_space.assignment }\OperatorTok{=}\NormalTok{ labels}
    \ControlFlowTok{for}\NormalTok{ k, centroid }\KeywordTok{in} \BuiltInTok{enumerate}\NormalTok{(centroids):}
\NormalTok{        mask }\OperatorTok{=}\NormalTok{ labels }\OperatorTok{==}\NormalTok{ k}
\NormalTok{        info }\OperatorTok{=}\NormalTok{ StateInfo(}
            \BuiltInTok{id}\OperatorTok{=}\NormalTok{k,}
\NormalTok{            label}\OperatorTok{=}\SpecialStringTok{f"state\_}\SpecialCharTok{\{}\NormalTok{k}\SpecialCharTok{\}}\SpecialStringTok{"}\NormalTok{,}
\NormalTok{            centroid}\OperatorTok{=}\NormalTok{centroid,}
\NormalTok{            radius}\OperatorTok{=}\BuiltInTok{float}\NormalTok{(np.mean([}\VariableTok{self}\NormalTok{.encoder.distance(centroid, Z[i])}
                                  \ControlFlowTok{for}\NormalTok{ i }\KeywordTok{in}\NormalTok{ np.where(mask)[}\DecValTok{0}\NormalTok{]])),}
\NormalTok{            count}\OperatorTok{=}\BuiltInTok{int}\NormalTok{(mask.}\BuiltInTok{sum}\NormalTok{()),}
\NormalTok{            proportion}\OperatorTok{=}\BuiltInTok{float}\NormalTok{(mask.mean()),}
\NormalTok{        )}
        \VariableTok{self}\NormalTok{.state\_space.add\_state(info)}

    \CommentTok{\# 4. Expert reliability}
    \VariableTok{self}\NormalTok{.expert\_reliability }\OperatorTok{=}\NormalTok{ ExpertReliability(...)}
\NormalTok{    result }\OperatorTok{=} \VariableTok{self}\NormalTok{.expert\_reliability.estimate(}
\NormalTok{        registry}\OperatorTok{=}\VariableTok{self}\NormalTok{.expert\_registry,}
\NormalTok{        X}\OperatorTok{=}\NormalTok{X, y}\OperatorTok{=}\NormalTok{y,}
\NormalTok{        state\_assignments}\OperatorTok{=}\NormalTok{labels,}
\NormalTok{        n\_states}\OperatorTok{=}\BuiltInTok{len}\NormalTok{(centroids),}
\NormalTok{    )}

    \CommentTok{\# 5. Data classification}
    \VariableTok{self}\NormalTok{.classifier }\OperatorTok{=}\NormalTok{ DataClassifier(config}\OperatorTok{=}\NormalTok{...)}
\NormalTok{    ...}

    \CommentTok{\# 6. Action policy}
\NormalTok{    ...}
\end{Highlighting}
\end{Shaded}

\begin{center}\rule{0.5\linewidth}{0.5pt}\end{center}

\subsection{7.
声明式领域配置}\label{ux58f0ux660eux5f0fux9886ux57dfux914dux7f6e}

每个领域只需一个 YAML 文件：

\subsubsection{7.1 MLIP 领域配置}\label{mlip-ux9886ux57dfux914dux7f6e}

\begin{Shaded}
\begin{Highlighting}[]
\CommentTok{\# domains/mlip.yaml}
\FunctionTok{domain}\KeywordTok{:}\AttributeTok{ mlip}
\FunctionTok{description}\KeywordTok{:}\AttributeTok{ }\StringTok{"机器学习原子间势函数的数据价值评估与专家路由"}

\FunctionTok{encoder}\KeywordTok{:}
\AttributeTok{  }\FunctionTok{class}\KeywordTok{:}\AttributeTok{ scx.encoders.mlip.ACEEncoder}
\AttributeTok{  }\FunctionTok{params}\KeywordTok{:}
\AttributeTok{    }\FunctionTok{descriptor}\KeywordTok{:}\AttributeTok{ soap}
\AttributeTok{    }\FunctionTok{rcut}\KeywordTok{:}\AttributeTok{ }\FloatTok{5.0}
\AttributeTok{    }\FunctionTok{nmax}\KeywordTok{:}\AttributeTok{ }\DecValTok{8}
\AttributeTok{    }\FunctionTok{lmax}\KeywordTok{:}\AttributeTok{ }\DecValTok{6}

\FunctionTok{experts}\KeywordTok{:}
\AttributeTok{  }\KeywordTok{{-}}\AttributeTok{ }\FunctionTok{name}\KeywordTok{:}\AttributeTok{ ACE}
\AttributeTok{    }\FunctionTok{type}\KeywordTok{:}\AttributeTok{ model}
\AttributeTok{    }\FunctionTok{predict\_fn}\KeywordTok{:}\AttributeTok{ }\StringTok{"path/to/ace\_model.ace"}
\AttributeTok{    }\FunctionTok{cost}\KeywordTok{:}\AttributeTok{ }\FloatTok{0.5}
\AttributeTok{  }\KeywordTok{{-}}\AttributeTok{ }\FunctionTok{name}\KeywordTok{:}\AttributeTok{ NEP}
\AttributeTok{    }\FunctionTok{type}\KeywordTok{:}\AttributeTok{ model}
\AttributeTok{    }\FunctionTok{predict\_fn}\KeywordTok{:}\AttributeTok{ }\StringTok{"path/to/nep\_model.txt"}
\AttributeTok{    }\FunctionTok{cost}\KeywordTok{:}\AttributeTok{ }\FloatTok{0.3}
\AttributeTok{  }\KeywordTok{{-}}\AttributeTok{ }\FunctionTok{name}\KeywordTok{:}\AttributeTok{ MACE}
\AttributeTok{    }\FunctionTok{type}\KeywordTok{:}\AttributeTok{ model}
\AttributeTok{    }\FunctionTok{predict\_fn}\KeywordTok{:}\AttributeTok{ }\StringTok{"path/to/mace\_model.pt"}
\AttributeTok{    }\FunctionTok{cost}\KeywordTok{:}\AttributeTok{ }\FloatTok{1.0}
\AttributeTok{  }\KeywordTok{{-}}\AttributeTok{ }\FunctionTok{name}\KeywordTok{:}\AttributeTok{ DFT}
\AttributeTok{    }\FunctionTok{type}\KeywordTok{:}\AttributeTok{ simulation}
\AttributeTok{    }\FunctionTok{command}\KeywordTok{:}\AttributeTok{ }\StringTok{"vasp\_std"}
\AttributeTok{    }\FunctionTok{cost}\KeywordTok{:}\AttributeTok{ }\FloatTok{100.0}

\FunctionTok{state\_config}\KeywordTok{:}
\AttributeTok{  }\FunctionTok{n\_states}\KeywordTok{:}\AttributeTok{ }\DecValTok{20}
\AttributeTok{  }\FunctionTok{cluster\_method}\KeywordTok{:}\AttributeTok{ hdbscan}
\AttributeTok{  }\FunctionTok{min\_cluster\_size}\KeywordTok{:}\AttributeTok{ }\DecValTok{10}

\FunctionTok{thresholds}\KeywordTok{:}
\AttributeTok{  }\FunctionTok{error\_high}\KeywordTok{:}\AttributeTok{ }\FloatTok{0.05}\CommentTok{       \# eV/atom}
\AttributeTok{  }\FunctionTok{density\_high}\KeywordTok{:}\AttributeTok{ }\FloatTok{0.05}
\AttributeTok{  }\FunctionTok{consistency\_high}\KeywordTok{:}\AttributeTok{ }\FloatTok{0.7}
\AttributeTok{  }\FunctionTok{redundancy\_high}\KeywordTok{:}\AttributeTok{ }\FloatTok{0.8}
\AttributeTok{  }\FunctionTok{noise\_high}\KeywordTok{:}\AttributeTok{ }\FloatTok{0.5}
\AttributeTok{  }\FunctionTok{expert\_gap}\KeywordTok{:}\AttributeTok{ }\FloatTok{0.3}

\FunctionTok{loss\_fn}\KeywordTok{:}\AttributeTok{ mse}\CommentTok{              \# 或 "mae", "rmse"}

\FunctionTok{reliability}\KeywordTok{:}
\AttributeTok{  }\FunctionTok{method}\KeywordTok{:}\AttributeTok{ supervised}
\AttributeTok{  }\FunctionTok{alpha}\KeywordTok{:}\AttributeTok{ }\FloatTok{1.0}
\AttributeTok{  }\FunctionTok{min\_samples}\KeywordTok{:}\AttributeTok{ }\DecValTok{5}

\FunctionTok{actions}\KeywordTok{:}
\AttributeTok{  }\KeywordTok{{-}}\AttributeTok{ acquire}
\AttributeTok{  }\KeywordTok{{-}}\AttributeTok{ compress}
\AttributeTok{  }\KeywordTok{{-}}\AttributeTok{ route}
\AttributeTok{  }\KeywordTok{{-}}\AttributeTok{ high\_fidelity}
\end{Highlighting}
\end{Shaded}

\subsubsection{7.2 Health
领域配置}\label{health-ux9886ux57dfux914dux7f6e}

\begin{Shaded}
\begin{Highlighting}[]
\CommentTok{\# domains/health.yaml}
\FunctionTok{domain}\KeywordTok{:}\AttributeTok{ health\_image}
\FunctionTok{description}\KeywordTok{:}\AttributeTok{ }\StringTok{"医学图像数据价值审计与专家复核调度"}

\FunctionTok{encoder}\KeywordTok{:}
\AttributeTok{  }\FunctionTok{class}\KeywordTok{:}\AttributeTok{ scx.encoders.vision.CNNEncoder}
\AttributeTok{  }\FunctionTok{params}\KeywordTok{:}
\AttributeTok{    }\FunctionTok{backbone}\KeywordTok{:}\AttributeTok{ resnet18}
\AttributeTok{    }\FunctionTok{device}\KeywordTok{:}\AttributeTok{ cuda}

\FunctionTok{experts}\KeywordTok{:}
\AttributeTok{  }\KeywordTok{{-}}\AttributeTok{ }\FunctionTok{name}\KeywordTok{:}\AttributeTok{ ai\_model\_v1}
\AttributeTok{    }\FunctionTok{type}\KeywordTok{:}\AttributeTok{ model}
\AttributeTok{    }\FunctionTok{checkpoint}\KeywordTok{:}\AttributeTok{ }\StringTok{"checkpoints/v1.pth"}
\AttributeTok{    }\FunctionTok{cost}\KeywordTok{:}\AttributeTok{ }\FloatTok{0.1}
\AttributeTok{  }\KeywordTok{{-}}\AttributeTok{ }\FunctionTok{name}\KeywordTok{:}\AttributeTok{ ai\_model\_v2}
\AttributeTok{    }\FunctionTok{type}\KeywordTok{:}\AttributeTok{ model}
\AttributeTok{    }\FunctionTok{checkpoint}\KeywordTok{:}\AttributeTok{ }\StringTok{"checkpoints/v2.pth"}
\AttributeTok{    }\FunctionTok{cost}\KeywordTok{:}\AttributeTok{ }\FloatTok{0.2}
\AttributeTok{  }\KeywordTok{{-}}\AttributeTok{ }\FunctionTok{name}\KeywordTok{:}\AttributeTok{ junior\_doctor}
\AttributeTok{    }\FunctionTok{type}\KeywordTok{:}\AttributeTok{ human}
\AttributeTok{    }\FunctionTok{cost}\KeywordTok{:}\AttributeTok{ }\FloatTok{50.0}
\AttributeTok{  }\KeywordTok{{-}}\AttributeTok{ }\FunctionTok{name}\KeywordTok{:}\AttributeTok{ senior\_doctor}
\AttributeTok{    }\FunctionTok{type}\KeywordTok{:}\AttributeTok{ human}
\AttributeTok{    }\FunctionTok{cost}\KeywordTok{:}\AttributeTok{ }\FloatTok{200.0}

\FunctionTok{state\_config}\KeywordTok{:}
\AttributeTok{  }\FunctionTok{n\_states}\KeywordTok{:}\AttributeTok{ }\DecValTok{30}
\AttributeTok{  }\FunctionTok{cluster\_method}\KeywordTok{:}\AttributeTok{ kmeans}
\AttributeTok{  }\FunctionTok{random\_state}\KeywordTok{:}\AttributeTok{ }\DecValTok{42}

\FunctionTok{thresholds}\KeywordTok{:}
\AttributeTok{  }\FunctionTok{error\_high}\KeywordTok{:}\AttributeTok{ }\FloatTok{0.3}
\AttributeTok{  }\FunctionTok{density\_high}\KeywordTok{:}\AttributeTok{ }\FloatTok{0.02}
\AttributeTok{  }\FunctionTok{consistency\_high}\KeywordTok{:}\AttributeTok{ }\FloatTok{0.6}
\AttributeTok{  }\FunctionTok{redundancy\_high}\KeywordTok{:}\AttributeTok{ }\FloatTok{0.85}
\AttributeTok{  }\FunctionTok{noise\_high}\KeywordTok{:}\AttributeTok{ }\FloatTok{0.4}
\AttributeTok{  }\FunctionTok{expert\_gap}\KeywordTok{:}\AttributeTok{ }\FloatTok{0.25}

\FunctionTok{loss\_fn}\KeywordTok{:}\AttributeTok{ cross\_entropy}

\FunctionTok{reliability}\KeywordTok{:}
\AttributeTok{  }\FunctionTok{method}\KeywordTok{:}\AttributeTok{ hybrid}
\AttributeTok{  }\FunctionTok{alpha}\KeywordTok{:}\AttributeTok{ }\FloatTok{2.0}

\FunctionTok{actions}\KeywordTok{:}
\AttributeTok{  }\KeywordTok{{-}}\AttributeTok{ keep}
\AttributeTok{  }\KeywordTok{{-}}\AttributeTok{ compress}
\AttributeTok{  }\KeywordTok{{-}}\AttributeTok{ route}
\AttributeTok{  }\KeywordTok{{-}}\AttributeTok{ relabel}
\end{Highlighting}
\end{Shaded}

\subsubsection{7.3 Robot 领域配置}\label{robot-ux9886ux57dfux914dux7f6e}

\begin{Shaded}
\begin{Highlighting}[]
\CommentTok{\# domains/robot.yaml}
\FunctionTok{domain}\KeywordTok{:}\AttributeTok{ robot\_trajectory}
\FunctionTok{description}\KeywordTok{:}\AttributeTok{ }\StringTok{"机器人示范轨迹价值评估与专家策略路由"}

\FunctionTok{encoder}\KeywordTok{:}
\AttributeTok{  }\FunctionTok{class}\KeywordTok{:}\AttributeTok{ scx.encoders.trajectory.TrajectoryEncoder}
\AttributeTok{  }\FunctionTok{params}\KeywordTok{:}
\AttributeTok{    }\FunctionTok{state\_dim}\KeywordTok{:}\AttributeTok{ }\DecValTok{128}
\AttributeTok{    }\FunctionTok{use\_temporal}\KeywordTok{:}\AttributeTok{ }\CharTok{true}

\FunctionTok{experts}\KeywordTok{:}
\AttributeTok{  }\KeywordTok{{-}}\AttributeTok{ }\FunctionTok{name}\KeywordTok{:}\AttributeTok{ scripted\_policy}
\AttributeTok{    }\FunctionTok{type}\KeywordTok{:}\AttributeTok{ rule\_based}
\AttributeTok{    }\FunctionTok{cost}\KeywordTok{:}\AttributeTok{ }\FloatTok{0.01}
\AttributeTok{  }\KeywordTok{{-}}\AttributeTok{ }\FunctionTok{name}\KeywordTok{:}\AttributeTok{ rl\_policy}
\AttributeTok{    }\FunctionTok{type}\KeywordTok{:}\AttributeTok{ model}
\AttributeTok{    }\FunctionTok{checkpoint}\KeywordTok{:}\AttributeTok{ }\StringTok{"checkpoints/rl.pt"}
\AttributeTok{    }\FunctionTok{cost}\KeywordTok{:}\AttributeTok{ }\FloatTok{0.5}
\AttributeTok{  }\KeywordTok{{-}}\AttributeTok{ }\FunctionTok{name}\KeywordTok{:}\AttributeTok{ mpc\_controller}
\AttributeTok{    }\FunctionTok{type}\KeywordTok{:}\AttributeTok{ simulation}
\AttributeTok{    }\FunctionTok{command}\KeywordTok{:}\AttributeTok{ }\StringTok{"mpc\_solver"}
\AttributeTok{    }\FunctionTok{cost}\KeywordTok{:}\AttributeTok{ }\FloatTok{10.0}
\AttributeTok{  }\KeywordTok{{-}}\AttributeTok{ }\FunctionTok{name}\KeywordTok{:}\AttributeTok{ human\_teleop}
\AttributeTok{    }\FunctionTok{type}\KeywordTok{:}\AttributeTok{ human}
\AttributeTok{    }\FunctionTok{cost}\KeywordTok{:}\AttributeTok{ }\FloatTok{500.0}
\AttributeTok{  }\KeywordTok{{-}}\AttributeTok{ }\FunctionTok{name}\KeywordTok{:}\AttributeTok{ vla\_model}
\AttributeTok{    }\FunctionTok{type}\KeywordTok{:}\AttributeTok{ model}
\AttributeTok{    }\FunctionTok{cost}\KeywordTok{:}\AttributeTok{ }\FloatTok{2.0}

\FunctionTok{state\_config}\KeywordTok{:}
\AttributeTok{  }\FunctionTok{n\_states}\KeywordTok{:}\AttributeTok{ }\DecValTok{20}
\AttributeTok{  }\FunctionTok{cluster\_method}\KeywordTok{:}\AttributeTok{ hdbscan}
\AttributeTok{  }\FunctionTok{min\_cluster\_size}\KeywordTok{:}\AttributeTok{ }\DecValTok{5}

\FunctionTok{thresholds}\KeywordTok{:}
\AttributeTok{  }\FunctionTok{error\_high}\KeywordTok{:}\AttributeTok{ }\FloatTok{0.15}
\AttributeTok{  }\FunctionTok{noise\_high}\KeywordTok{:}\AttributeTok{ }\FloatTok{0.5}
\AttributeTok{  }\FunctionTok{redundancy\_high}\KeywordTok{:}\AttributeTok{ }\FloatTok{0.75}
\AttributeTok{  }\FunctionTok{expert\_gap}\KeywordTok{:}\AttributeTok{ }\FloatTok{0.2}

\FunctionTok{reliability}\KeywordTok{:}
\AttributeTok{  }\FunctionTok{method}\KeywordTok{:}\AttributeTok{ bayesian}
\AttributeTok{  }\FunctionTok{min\_samples}\KeywordTok{:}\AttributeTok{ }\DecValTok{3}

\FunctionTok{actions}\KeywordTok{:}
\AttributeTok{  }\KeywordTok{{-}}\AttributeTok{ keep}
\AttributeTok{  }\KeywordTok{{-}}\AttributeTok{ compress}
\AttributeTok{  }\KeywordTok{{-}}\AttributeTok{ route}
\AttributeTok{  }\KeywordTok{{-}}\AttributeTok{ high\_fidelity}\CommentTok{        \# 上真实机器人}
\AttributeTok{  }\KeywordTok{{-}}\AttributeTok{ acquire}
\end{Highlighting}
\end{Shaded}

\subsubsection{7.4 FEM/Simulation
领域配置}\label{femsimulation-ux9886ux57dfux914dux7f6e}

\begin{Shaded}
\begin{Highlighting}[]
\CommentTok{\# domains/fem.yaml}
\FunctionTok{domain}\KeywordTok{:}\AttributeTok{ fem\_simulation}
\FunctionTok{description}\KeywordTok{:}\AttributeTok{ }\StringTok{"有限元仿真多保真调度与数据压缩"}

\FunctionTok{encoder}\KeywordTok{:}
\AttributeTok{  }\FunctionTok{class}\KeywordTok{:}\AttributeTok{ scx.encoders.simulation.FEMEncoder}
\AttributeTok{  }\FunctionTok{params}\KeywordTok{:}
\AttributeTok{    }\FunctionTok{mesh\_encoder}\KeywordTok{:}\AttributeTok{ }\StringTok{"pointnet"}
\AttributeTok{    }\FunctionTok{response\_dim}\KeywordTok{:}\AttributeTok{ }\DecValTok{64}

\FunctionTok{experts}\KeywordTok{:}
\AttributeTok{  }\KeywordTok{{-}}\AttributeTok{ }\FunctionTok{name}\KeywordTok{:}\AttributeTok{ coarse\_mesh}
\AttributeTok{    }\FunctionTok{type}\KeywordTok{:}\AttributeTok{ simulation}
\AttributeTok{    }\FunctionTok{config}\KeywordTok{:}\AttributeTok{ }\StringTok{"configs/coarse.json"}
\AttributeTok{    }\FunctionTok{cost}\KeywordTok{:}\AttributeTok{ }\FloatTok{1.0}
\AttributeTok{  }\KeywordTok{{-}}\AttributeTok{ }\FunctionTok{name}\KeywordTok{:}\AttributeTok{ fine\_mesh}
\AttributeTok{    }\FunctionTok{type}\KeywordTok{:}\AttributeTok{ simulation}
\AttributeTok{    }\FunctionTok{config}\KeywordTok{:}\AttributeTok{ }\StringTok{"configs/fine.json"}
\AttributeTok{    }\FunctionTok{cost}\KeywordTok{:}\AttributeTok{ }\FloatTok{50.0}
\AttributeTok{  }\KeywordTok{{-}}\AttributeTok{ }\FunctionTok{name}\KeywordTok{:}\AttributeTok{ neural\_operator}
\AttributeTok{    }\FunctionTok{type}\KeywordTok{:}\AttributeTok{ model}
\AttributeTok{    }\FunctionTok{checkpoint}\KeywordTok{:}\AttributeTok{ }\StringTok{"checkpoints/don.pt"}
\AttributeTok{    }\FunctionTok{cost}\KeywordTok{:}\AttributeTok{ }\FloatTok{0.1}
\AttributeTok{  }\KeywordTok{{-}}\AttributeTok{ }\FunctionTok{name}\KeywordTok{:}\AttributeTok{ analytical}
\AttributeTok{    }\FunctionTok{type}\KeywordTok{:}\AttributeTok{ rule\_based}
\AttributeTok{    }\FunctionTok{cost}\KeywordTok{:}\AttributeTok{ }\FloatTok{0.01}

\FunctionTok{state\_config}\KeywordTok{:}
\AttributeTok{  }\FunctionTok{n\_states}\KeywordTok{:}\AttributeTok{ }\DecValTok{15}
\AttributeTok{  }\FunctionTok{cluster\_method}\KeywordTok{:}\AttributeTok{ spectral}

\FunctionTok{thresholds}\KeywordTok{:}
\AttributeTok{  }\FunctionTok{error\_high}\KeywordTok{:}\AttributeTok{ }\FloatTok{0.1}
\AttributeTok{  }\FunctionTok{redundancy\_high}\KeywordTok{:}\AttributeTok{ }\FloatTok{0.8}
\AttributeTok{  }\FunctionTok{noise\_high}\KeywordTok{:}\AttributeTok{ }\FloatTok{0.4}
\AttributeTok{  }\FunctionTok{expert\_gap}\KeywordTok{:}\AttributeTok{ }\FloatTok{0.3}

\FunctionTok{actions}\KeywordTok{:}
\AttributeTok{  }\KeywordTok{{-}}\AttributeTok{ keep}
\AttributeTok{  }\KeywordTok{{-}}\AttributeTok{ compress}
\AttributeTok{  }\KeywordTok{{-}}\AttributeTok{ route}
\AttributeTok{  }\KeywordTok{{-}}\AttributeTok{ high\_fidelity}
\end{Highlighting}
\end{Shaded}

\subsubsection{7.5 Tabular/Benchmark
配置}\label{tabularbenchmark-ux914dux7f6e}

\begin{Shaded}
\begin{Highlighting}[]
\CommentTok{\# domains/cifar.yaml}
\FunctionTok{domain}\KeywordTok{:}\AttributeTok{ cifar\_image}
\FunctionTok{description}\KeywordTok{:}\AttributeTok{ }\StringTok{"CIFAR 图像数据价值评估"}

\FunctionTok{encoder}\KeywordTok{:}
\AttributeTok{  }\FunctionTok{class}\KeywordTok{:}\AttributeTok{ scx.encoders.vision.CNNEncoder}
\AttributeTok{  }\FunctionTok{params}\KeywordTok{:}
\AttributeTok{    }\FunctionTok{backbone}\KeywordTok{:}\AttributeTok{ simple\_cnn}
\AttributeTok{    }\FunctionTok{device}\KeywordTok{:}\AttributeTok{ cpu}

\FunctionTok{experts}\KeywordTok{:}
\AttributeTok{  }\KeywordTok{{-}}\AttributeTok{ }\FunctionTok{name}\KeywordTok{:}\AttributeTok{ cnn\_small}
\AttributeTok{    }\FunctionTok{type}\KeywordTok{:}\AttributeTok{ model}
\AttributeTok{    }\FunctionTok{cost}\KeywordTok{:}\AttributeTok{ }\FloatTok{0.1}
\AttributeTok{  }\KeywordTok{{-}}\AttributeTok{ }\FunctionTok{name}\KeywordTok{:}\AttributeTok{ cnn\_large}
\AttributeTok{    }\FunctionTok{type}\KeywordTok{:}\AttributeTok{ model}
\AttributeTok{    }\FunctionTok{cost}\KeywordTok{:}\AttributeTok{ }\FloatTok{0.5}
\AttributeTok{  }\KeywordTok{{-}}\AttributeTok{ }\FunctionTok{name}\KeywordTok{:}\AttributeTok{ resnet50}
\AttributeTok{    }\FunctionTok{type}\KeywordTok{:}\AttributeTok{ model}
\AttributeTok{    }\FunctionTok{cost}\KeywordTok{:}\AttributeTok{ }\FloatTok{1.0}

\FunctionTok{state\_config}\KeywordTok{:}
\AttributeTok{  }\FunctionTok{n\_states}\KeywordTok{:}\AttributeTok{ }\DecValTok{25}
\AttributeTok{  }\FunctionTok{cluster\_method}\KeywordTok{:}\AttributeTok{ kmeans}
\AttributeTok{  }\FunctionTok{random\_state}\KeywordTok{:}\AttributeTok{ }\DecValTok{42}

\FunctionTok{thresholds}\KeywordTok{:}
\AttributeTok{  }\FunctionTok{error\_high}\KeywordTok{:}\AttributeTok{ }\FloatTok{0.3}
\AttributeTok{  }\FunctionTok{density\_high}\KeywordTok{:}\AttributeTok{ }\FloatTok{0.04}
\AttributeTok{  }\FunctionTok{consistency\_high}\KeywordTok{:}\AttributeTok{ }\FloatTok{0.6}
\AttributeTok{  }\FunctionTok{redundancy\_high}\KeywordTok{:}\AttributeTok{ }\FloatTok{0.85}
\AttributeTok{  }\FunctionTok{noise\_high}\KeywordTok{:}\AttributeTok{ }\FloatTok{0.5}
\AttributeTok{  }\FunctionTok{expert\_gap}\KeywordTok{:}\AttributeTok{ }\FloatTok{0.2}

\FunctionTok{loss\_fn}\KeywordTok{:}\AttributeTok{ cross\_entropy}

\FunctionTok{actions}\KeywordTok{:}
\AttributeTok{  }\KeywordTok{{-}}\AttributeTok{ keep}
\AttributeTok{  }\KeywordTok{{-}}\AttributeTok{ compress}
\AttributeTok{  }\KeywordTok{{-}}\AttributeTok{ route}
\AttributeTok{  }\KeywordTok{{-}}\AttributeTok{ relabel}
\end{Highlighting}
\end{Shaded}

\begin{center}\rule{0.5\linewidth}{0.5pt}\end{center}

\subsection{8. 文件夹结构}\label{ux6587ux4ef6ux5939ux7ed3ux6784}

\begin{verbatim}
scx/
├── core/                      # 🟢 通用核心（已有，基本不变）
│   ├── __init__.py
│   ├── framework.py           # SCXFramework（重构为领域引擎调度）
│   ├── config.py              # SCXConfig（扩展领域配置加载）
│   ├── metrics.py             # SCXMetrics（不变）
│   ├── online.py              # OnlineSCXFramework（不变）
│   ├── domain_engine.py       # 🆕 SCXDomainEngine（新入口）
│   ├── adapter.py             # 🆕 自动推导 Adapter
│   ├── audit.py               # 🆕 SCXAuditReport
│   ├── expert.py              # 🆕 SCXExpert 抽象
│   └── datapoint.py           # 🆕 SCXDataPoint 抽象
│
├── state/                     # 🟢 状态管理（已有，不变）
│   ├── space.py               # StateSpace + StateInfo
│   ├── discovery.py           # 状态发现（不变）
│   ├── assignment.py          # 状态分配（不变）
│   └── robustess.py           # 鲁棒性分析（不变）
│
├── expert/                    # 🟢 专家管理（已有，不变）
│   ├── registry.py            # ExpertRegistry（不变）
│   ├── reliability.py         # ExpertReliability（不变）
│   ├── router.py              # ExpertRouter（不变）
│   └── conflict.py            # ExpertConflict（不变）
│
├── valuation/                 # 🟢 数据估值（已有，不变）
│   ├── classifier.py          # DataClassifier（不变）
│   ├── state_value.py         # 状态价值（不变）
│   ├── redundancy.py          # 冗余度（不变）
│   ├── noise_score.py         # 噪声评分（不变）
│   ├── learnability.py        # 可学习性（不变）
│   ├── influence.py           # Influence 分数（不变）
│   └── adaptive.py            # 自适应估值（不变）
│
├── action/                    # 🟢 动作策略（已有，不变）
│   ├── policy.py              # ActionPolicy（不变）
│   ├── compress.py            # 压缩策略（不变）
│   └── acquisition.py         # 获取策略（不变）
│
├── encoders/                  # 🆕 领域编码器（统一接口）
│   ├── __init__.py
│   ├── base.py                # SCXStateEncoder 抽象基类
│   ├── mlip.py                # ACE/SOAP descriptor encoder
│   ├── vision.py              # CNN/ViT/CLIP encoder
│   ├── tabular.py             # 表格数据 encoder
│   ├── trajectory.py          # 机器人轨迹 encoder
│   ├── simulation.py          # FEM/PDE 仿真 encoder
│   ├── text.py                # LLM/NLP embedding encoder
│   └── health.py              # 医学图像 encoder (从 scx-health 迁移)
│
├── domains/                   # 🆕 领域配置（声明式）
│   ├── __init__.py
│   ├── mlip.yaml
│   ├── health.yaml
│   ├── cifar.yaml
│   ├── robot.yaml
│   ├── fem.yaml
│   ├── tabular.yaml
│   ├── text.yaml
│   └── registry.py            # 🆕 领域注册表
│
├── plugins/                   # 🆕 插件（可选扩展）
│   ├── __init__.py
│   ├── compress.py            # SCX-Compress 插件
│   ├── online.py              # Online SCX 插件
│   └── influence.py           # Influence 插件
│
├── utils/                     # 🟢 工具（已有，扩展）
│   ├── helpers.py             # 辅助函数
│   ├── data_loader.py         # 数据加载
│   ├── visualization.py       # 可视化
│   └── evaluation.py          # 评估工具
│
└── __init__.py
\end{verbatim}

\begin{center}\rule{0.5\linewidth}{0.5pt}\end{center}

\subsection{9. SCX Domain
注册与生命周期}\label{scx-domain-ux6ce8ux518cux4e0eux751fux547dux5468ux671f}

\subsubsection{9.1 领域注册表}\label{ux9886ux57dfux6ce8ux518cux8868}

\begin{Shaded}
\begin{Highlighting}[]
\CommentTok{\# scx/domains/registry.py}

\KeywordTok{class}\NormalTok{ DomainRegistry:}
    \CommentTok{"""全局领域配置注册表，支持按名称查找。"""}

\NormalTok{    \_domains: }\BuiltInTok{dict}\NormalTok{[}\BuiltInTok{str}\NormalTok{, }\BuiltInTok{dict}\NormalTok{] }\OperatorTok{=}\NormalTok{ \{\}}

    \AttributeTok{@classmethod}
    \KeywordTok{def}\NormalTok{ register(cls, name: }\BuiltInTok{str}\NormalTok{, config: }\BuiltInTok{dict} \OperatorTok{|} \BuiltInTok{str}\NormalTok{):}
        \CommentTok{"""注册一个领域。config 可以是 dict 或 YAML 文件路径。"""}
        \ControlFlowTok{if} \BuiltInTok{isinstance}\NormalTok{(config, }\BuiltInTok{str}\NormalTok{):}
            \ImportTok{import}\NormalTok{ yaml}
            \ControlFlowTok{with} \BuiltInTok{open}\NormalTok{(config) }\ImportTok{as}\NormalTok{ f:}
\NormalTok{                config }\OperatorTok{=}\NormalTok{ yaml.safe\_load(f)}
\NormalTok{        cls.\_domains[name] }\OperatorTok{=}\NormalTok{ config}

    \AttributeTok{@classmethod}
    \KeywordTok{def}\NormalTok{ get(cls, name: }\BuiltInTok{str}\NormalTok{) }\OperatorTok{{-}\textgreater{}} \BuiltInTok{dict}\NormalTok{:}
        \ControlFlowTok{if}\NormalTok{ name }\KeywordTok{not} \KeywordTok{in}\NormalTok{ cls.\_domains:}
            \ControlFlowTok{raise} \PreprocessorTok{KeyError}\NormalTok{(}\SpecialStringTok{f"Domain \textquotesingle{}}\SpecialCharTok{\{}\NormalTok{name}\SpecialCharTok{\}}\SpecialStringTok{\textquotesingle{} not registered. "}
                           \SpecialStringTok{f"Available: }\SpecialCharTok{\{}\BuiltInTok{list}\NormalTok{(cls.\_domains.keys())}\SpecialCharTok{\}}\SpecialStringTok{"}\NormalTok{)}
        \ControlFlowTok{return}\NormalTok{ cls.\_domains[name]}

    \AttributeTok{@classmethod}
    \KeywordTok{def}\NormalTok{ list\_domains(cls) }\OperatorTok{{-}\textgreater{}} \BuiltInTok{list}\NormalTok{[}\BuiltInTok{str}\NormalTok{]:}
        \ControlFlowTok{return} \BuiltInTok{list}\NormalTok{(cls.\_domains.keys())}

    \AttributeTok{@classmethod}
    \KeywordTok{def}\NormalTok{ load\_all(cls, domains\_dir: }\BuiltInTok{str} \OperatorTok{=} \VariableTok{None}\NormalTok{):}
        \CommentTok{"""从目录加载所有 YAML 领域配置。"""}
        \ControlFlowTok{if}\NormalTok{ domains\_dir }\KeywordTok{is} \VariableTok{None}\NormalTok{:}
\NormalTok{            domains\_dir }\OperatorTok{=}\NormalTok{ Path(}\VariableTok{\_\_file\_\_}\NormalTok{).parent}
        \ImportTok{import}\NormalTok{ glob}
        \ControlFlowTok{for}\NormalTok{ yaml\_file }\KeywordTok{in}\NormalTok{ glob.glob(}\BuiltInTok{str}\NormalTok{(domains\_dir }\OperatorTok{/} \StringTok{"*.yaml"}\NormalTok{)):}
\NormalTok{            name }\OperatorTok{=}\NormalTok{ Path(yaml\_file).stem}
\NormalTok{            cls.register(name, yaml\_file)}

\CommentTok{\# 启动时自动加载所有领域配置}
\NormalTok{DomainRegistry.load\_all()}
\end{Highlighting}
\end{Shaded}

\subsubsection{9.2 生命周期}\label{ux751fux547dux5468ux671f}

\begin{verbatim}
初始化:
  1. DomainRegistry.get("mlip") → 加载 YAML
  2. SCXDomainEngine(domain_config) → 构建引擎
  3. 自动实例化 Encoder
  4. 自动创建 ExpertRegistry
  5. 自动配置阈值

运行:
  1. engine.fit(X, y) → 完整 pipeline
  2. engine.audit() → 输出审计报告
  3. engine.value(x) → 单样本价值
  4. engine.recommend_action(state_id)

持久化:
  1. engine.save("my_project.scx") → 序列化全部状态
  2. SCXDomainEngine.load("my_project.scx") → 恢复

在线模式:
  1. OnlineSCXFramework(initial_centroids, M) → 流式处理
  2. process_sample(x, expert_id, loss) → 增量更新
\end{verbatim}

\begin{center}\rule{0.5\linewidth}{0.5pt}\end{center}

\subsection{10. 插件系统}\label{ux63d2ux4ef6ux7cfbux7edf}

插件是可选的扩展模块，提供超越 Core 的额外功能。

\subsubsection{10.1 插件接口}\label{ux63d2ux4ef6ux63a5ux53e3}

\begin{Shaded}
\begin{Highlighting}[]
\CommentTok{\# scx/plugins/base.py}

\KeywordTok{class}\NormalTok{ SCXPlugin(ABC):}
    \CommentTok{"""SCX 插件基类。"""}

    \AttributeTok{@abstractmethod}
    \KeywordTok{def}\NormalTok{ name(}\VariableTok{self}\NormalTok{) }\OperatorTok{{-}\textgreater{}} \BuiltInTok{str}\NormalTok{: ...}

    \AttributeTok{@abstractmethod}
    \KeywordTok{def}\NormalTok{ attach(}\VariableTok{self}\NormalTok{, engine: SCXDomainEngine) }\OperatorTok{{-}\textgreater{}} \VariableTok{None}\NormalTok{:}
        \CommentTok{"""将插件挂载到领域引擎上。"""}
\NormalTok{        ...}

    \KeywordTok{def}\NormalTok{ on\_fit\_complete(}\VariableTok{self}\NormalTok{, engine: SCXDomainEngine) }\OperatorTok{{-}\textgreater{}} \VariableTok{None}\NormalTok{:}
        \CommentTok{"""fit() 完成后回调。"""}
\NormalTok{        ...}

    \KeywordTok{def}\NormalTok{ on\_classify\_complete(}\VariableTok{self}\NormalTok{, engine: SCXDomainEngine) }\OperatorTok{{-}\textgreater{}} \VariableTok{None}\NormalTok{:}
        \CommentTok{"""分类完成后回调。"""}
\NormalTok{        ...}
\end{Highlighting}
\end{Shaded}

\subsubsection{10.2 内置插件}\label{ux5185ux7f6eux63d2ux4ef6}

\begin{longtable}[]{@{}
  >{\raggedright\arraybackslash}p{(\linewidth - 4\tabcolsep) * \real{0.3333}}
  >{\raggedright\arraybackslash}p{(\linewidth - 4\tabcolsep) * \real{0.3333}}
  >{\raggedright\arraybackslash}p{(\linewidth - 4\tabcolsep) * \real{0.3333}}@{}}
\toprule\noalign{}
\begin{minipage}[b]{\linewidth}\raggedright
插件
\end{minipage} & \begin{minipage}[b]{\linewidth}\raggedright
文件
\end{minipage} & \begin{minipage}[b]{\linewidth}\raggedright
功能
\end{minipage} \\
\midrule\noalign{}
\endhead
\bottomrule\noalign{}
\endlastfoot
\texttt{CompressPlugin} & \texttt{plugins/compress.py} &
数据压缩策略（corest selection） \\
\texttt{OnlinePlugin} & \texttt{plugins/online.py} & 流式在线监测 \\
\texttt{InfluencePlugin} & \texttt{plugins/influence.py} & Influence
function 数据估值 \\
\end{longtable}

\subsubsection{10.3 插件的使用}\label{ux63d2ux4ef6ux7684ux4f7fux7528}

\begin{Shaded}
\begin{Highlighting}[]
\CommentTok{\# 通过配置启用插件}
\NormalTok{engine }\OperatorTok{=}\NormalTok{ SCXDomainEngine(}\StringTok{"domains/mlip.yaml"}\NormalTok{)}
\NormalTok{engine.attach\_plugin(OnlinePlugin(tau}\OperatorTok{=}\FloatTok{0.5}\NormalTok{, decay}\OperatorTok{=}\FloatTok{0.95}\NormalTok{))}
\NormalTok{engine.fit(X, y)}
\end{Highlighting}
\end{Shaded}

\begin{center}\rule{0.5\linewidth}{0.5pt}\end{center}

\subsection{11.
与现有代码的迁移路径}\label{ux4e0eux73b0ux6709ux4ee3ux7801ux7684ux8fc1ux79fbux8defux5f84}

\subsubsection{11.1 迁移原则}\label{ux8fc1ux79fbux539fux5219}

\begin{enumerate}
\def\labelenumi{\arabic{enumi}.}
\tightlist
\item
  \textbf{不破坏现有代码}：已有的 \texttt{scx.core}, \texttt{scx.state},
  \texttt{scx.expert}, \texttt{scx.valuation}, \texttt{scx.action} 不变
\item
  \textbf{增量重构}：先加新结构，再逐步迁移旧的 domain 代码
\item
  \textbf{向后兼容}：\texttt{SCXFramework} 保留，但标记为
  deprecated，推荐使用 \texttt{SCXDomainEngine}
\end{enumerate}

\subsubsection{11.2 迁移步骤}\label{ux8fc1ux79fbux6b65ux9aa4}

\paragraph{Phase
1：新增结构（无破坏）}\label{phase-1ux65b0ux589eux7ed3ux6784ux65e0ux7834ux574f}

\begin{verbatim}
Step 1: 创建 scx/encoders/base.py          # SCXStateEncoder ABC
Step 2: 创建 scx/encoders/mlip.py          # MLIP Encoder
Step 3: 创建 scx/encoders/vision.py        # Vision Encoder
Step 4: 创建 scx/core/domain_engine.py     # SCXDomainEngine
Step 5: 创建 scx/core/adapter.py           # 自动推导逻辑
Step 6: 创建 scx/core/audit.py             # 审计报告
Step 7: 创建 scx/domains/registry.py       # 领域注册表
Step 8: 创建配置文件 mlip.yaml, health.yaml, cifar.yaml
\end{verbatim}

\paragraph{Phase 2：迁移
scx-health}\label{phase-2ux8fc1ux79fb-scx-health}

\begin{verbatim}
当前: scx-health/src/scx_health/encoder.py  (独立包)
目标: scx/encoders/health.py                (统一接口)

迁移:
1. 将 SimpleCNN, ResNetEncoder 移植到 scx/encoders/health.py
2. 实现 SCXStateEncoder 接口（encode/distance/cluster）
3. 创建 scx/domains/health.yaml
4. 验证原 scx-health 实验仍能运行
5. 原 scx-health 包保留为 legacy，标记 deprecated
\end{verbatim}

\paragraph{Phase 3：迁移 CIFAR
实验}\label{phase-3ux8fc1ux79fb-cifar-ux5b9eux9a8c}

\begin{verbatim}
当前: experiments/cifar/utils.py  (独立 encoder)
目标: scx/encoders/vision.py      (复用统一接口)

迁移:
1. CIFAR 的 SimpleCNN 可以用 scx/encoders/vision.CNNEncoder
2. 创建 scx/domains/cifar.yaml
3. 实验脚本改为使用 SCXDomainEngine
4. 原 experiments/cifar/ 保留，标记为 legacy
\end{verbatim}

\paragraph{Phase
4：新领域接入}\label{phase-4ux65b0ux9886ux57dfux63a5ux5165}

\begin{verbatim}
机器人:  scx/encoders/trajectory.py + domains/robot.yaml
FEM:     scx/encoders/simulation.py + domains/fem.yaml
表格:    scx/encoders/tabular.py + domains/tabular.yaml
文本:    scx/encoders/text.py + domains/text.yaml
\end{verbatim}

\subsubsection{11.3
新老接口对比}\label{ux65b0ux8001ux63a5ux53e3ux5bf9ux6bd4}

\begin{longtable}[]{@{}
  >{\raggedright\arraybackslash}p{(\linewidth - 4\tabcolsep) * \real{0.2727}}
  >{\raggedright\arraybackslash}p{(\linewidth - 4\tabcolsep) * \real{0.3636}}
  >{\raggedright\arraybackslash}p{(\linewidth - 4\tabcolsep) * \real{0.3636}}@{}}
\toprule\noalign{}
\begin{minipage}[b]{\linewidth}\raggedright
场景
\end{minipage} & \begin{minipage}[b]{\linewidth}\raggedright
旧方式
\end{minipage} & \begin{minipage}[b]{\linewidth}\raggedright
新方式
\end{minipage} \\
\midrule\noalign{}
\endhead
\bottomrule\noalign{}
\endlastfoot
创建框架 & \texttt{SCXFramework(config)} &
\texttt{SCXDomainEngine("domains/mlip.yaml")} \\
添加领域 & 写 adapter 类，修改 framework & 写 Encoder +
YAML，直接注册 \\
配置阈值 & 硬编码在 config.py & YAML 中声明 \\
数据编码 & 调用方自己处理 & Encoder.encode() 统一接口 \\
聚类 & 调用方自己选择算法 & Encoder.cluster() 统一接口 \\
特征 & 手动传递 phi(X) & encoder.encode\_batch(X) 自动 \\
审计报告 & 不统一 & SCXAuditReport 标准格式 \\
领域管理 & 无 & DomainRegistry \\
\end{longtable}

\begin{center}\rule{0.5\linewidth}{0.5pt}\end{center}

\subsection{12.
新增领域的步骤}\label{ux65b0ux589eux9886ux57dfux7684ux6b65ux9aa4}

\subsubsection{12.1
标准三步流程}\label{ux6807ux51c6ux4e09ux6b65ux6d41ux7a0b}

\begin{verbatim}
第 1 步 ── 写 Encoder
第 2 步 ── 写 YAML
第 3 步 ── 跑测试
\end{verbatim}

\subsubsection{12.2 详细步骤}\label{ux8be6ux7ec6ux6b65ux9aa4}

\paragraph{第 1 步：写 Encoder（约 30-100
行）}\label{ux7b2c-1-ux6b65ux5199-encoderux7ea6-30-100-ux884c}

\begin{Shaded}
\begin{Highlighting}[]
\CommentTok{\# scx/encoders/my\_domain.py}

\ImportTok{from}\NormalTok{ scx.encoders.base }\ImportTok{import}\NormalTok{ SCXStateEncoder}
\ImportTok{import}\NormalTok{ numpy }\ImportTok{as}\NormalTok{ np}

\KeywordTok{class}\NormalTok{ MyDomainEncoder(SCXStateEncoder):}
    \CommentTok{"""我的领域编码器。"""}

    \KeywordTok{def} \FunctionTok{\_\_init\_\_}\NormalTok{(}\VariableTok{self}\NormalTok{, param1: }\BuiltInTok{int} \OperatorTok{=} \DecValTok{10}\NormalTok{, param2: }\BuiltInTok{str} \OperatorTok{=} \StringTok{"default"}\NormalTok{):}
        \VariableTok{self}\NormalTok{.param1 }\OperatorTok{=}\NormalTok{ param1}
        \VariableTok{self}\NormalTok{.param2 }\OperatorTok{=}\NormalTok{ param2}
        \VariableTok{self}\NormalTok{.\_feature\_dim }\OperatorTok{=}\NormalTok{ param1}

    \KeywordTok{def}\NormalTok{ encode(}\VariableTok{self}\NormalTok{, x: Any) }\OperatorTok{{-}\textgreater{}}\NormalTok{ np.ndarray:}
        \CommentTok{\# 1. 预处理 x}
        \CommentTok{\# 2. 提取特征向量}
        \CommentTok{\# 3. 返回 (d,) 的 np.ndarray}
\NormalTok{        ...}

    \KeywordTok{def}\NormalTok{ distance(}\VariableTok{self}\NormalTok{, a: np.ndarray, b: np.ndarray) }\OperatorTok{{-}\textgreater{}} \BuiltInTok{float}\NormalTok{:}
        \CommentTok{\# 自定义距离度量}
        \ControlFlowTok{return} \BuiltInTok{float}\NormalTok{(np.linalg.norm(a }\OperatorTok{{-}}\NormalTok{ b))}

    \KeywordTok{def}\NormalTok{ cluster(}\VariableTok{self}\NormalTok{, Z: np.ndarray, }\OperatorTok{**}\NormalTok{kwargs) }\OperatorTok{{-}\textgreater{}} \BuiltInTok{tuple}\NormalTok{[np.ndarray, np.ndarray]:}
        \CommentTok{\# 选择适合本领域的聚类方法}
        \ImportTok{from}\NormalTok{ sklearn.cluster }\ImportTok{import}\NormalTok{ KMeans}
\NormalTok{        n }\OperatorTok{=}\NormalTok{ kwargs.get(}\StringTok{"n\_states"}\NormalTok{, }\DecValTok{10}\NormalTok{)}
\NormalTok{        km }\OperatorTok{=}\NormalTok{ KMeans(n\_clusters}\OperatorTok{=}\NormalTok{n, random\_state}\OperatorTok{=}\DecValTok{42}\NormalTok{)}
\NormalTok{        labels }\OperatorTok{=}\NormalTok{ km.fit\_predict(Z)}
        \ControlFlowTok{return}\NormalTok{ labels, km.cluster\_centers\_}
\end{Highlighting}
\end{Shaded}

\paragraph{第 2 步：写 YAML（约 20-40
行）}\label{ux7b2c-2-ux6b65ux5199-yamlux7ea6-20-40-ux884c}

\begin{Shaded}
\begin{Highlighting}[]
\CommentTok{\# scx/domains/my\_domain.yaml}
\FunctionTok{domain}\KeywordTok{:}\AttributeTok{ my\_domain}
\FunctionTok{description}\KeywordTok{:}\AttributeTok{ }\StringTok{"我的领域数据价值评估"}

\FunctionTok{encoder}\KeywordTok{:}
\AttributeTok{  }\FunctionTok{class}\KeywordTok{:}\AttributeTok{ scx.encoders.my\_domain.MyDomainEncoder}
\AttributeTok{  }\FunctionTok{params}\KeywordTok{:}
\AttributeTok{    }\FunctionTok{param1}\KeywordTok{:}\AttributeTok{ }\DecValTok{10}
\AttributeTok{    }\FunctionTok{param2}\KeywordTok{:}\AttributeTok{ custom}

\FunctionTok{experts}\KeywordTok{:}
\AttributeTok{  }\KeywordTok{{-}}\AttributeTok{ }\FunctionTok{name}\KeywordTok{:}\AttributeTok{ expert\_a}
\AttributeTok{    }\FunctionTok{type}\KeywordTok{:}\AttributeTok{ model}
\AttributeTok{    }\FunctionTok{cost}\KeywordTok{:}\AttributeTok{ }\FloatTok{1.0}
\AttributeTok{  }\KeywordTok{{-}}\AttributeTok{ }\FunctionTok{name}\KeywordTok{:}\AttributeTok{ expert\_b}
\AttributeTok{    }\FunctionTok{type}\KeywordTok{:}\AttributeTok{ model}
\AttributeTok{    }\FunctionTok{cost}\KeywordTok{:}\AttributeTok{ }\FloatTok{5.0}

\FunctionTok{state\_config}\KeywordTok{:}
\AttributeTok{  }\FunctionTok{n\_states}\KeywordTok{:}\AttributeTok{ }\DecValTok{15}
\AttributeTok{  }\FunctionTok{cluster\_method}\KeywordTok{:}\AttributeTok{ kmeans}

\FunctionTok{thresholds}\KeywordTok{:}
\AttributeTok{  }\FunctionTok{error\_high}\KeywordTok{:}\AttributeTok{ }\FloatTok{0.1}
\AttributeTok{  }\FunctionTok{noise\_high}\KeywordTok{:}\AttributeTok{ }\FloatTok{0.4}
\AttributeTok{  }\FunctionTok{redundancy\_high}\KeywordTok{:}\AttributeTok{ }\FloatTok{0.8}
\AttributeTok{  }\FunctionTok{expert\_gap}\KeywordTok{:}\AttributeTok{ }\FloatTok{0.25}

\FunctionTok{actions}\KeywordTok{:}
\AttributeTok{  }\KeywordTok{{-}}\AttributeTok{ keep}
\AttributeTok{  }\KeywordTok{{-}}\AttributeTok{ compress}
\AttributeTok{  }\KeywordTok{{-}}\AttributeTok{ route}
\end{Highlighting}
\end{Shaded}

\paragraph{第 3 步：跑测试}\label{ux7b2c-3-ux6b65ux8dd1ux6d4bux8bd5}

\begin{Shaded}
\begin{Highlighting}[]
\CommentTok{\# 测试脚本}

\ImportTok{from}\NormalTok{ scx.core.domain\_engine }\ImportTok{import}\NormalTok{ SCXDomainEngine}
\ImportTok{from}\NormalTok{ scx.domains.registry }\ImportTok{import}\NormalTok{ DomainRegistry}

\CommentTok{\# 1. 注册领域（自动或手动）}
\NormalTok{DomainRegistry.register(}\StringTok{"my\_domain"}\NormalTok{, }\StringTok{"scx/domains/my\_domain.yaml"}\NormalTok{)}

\CommentTok{\# 2. 创建引擎}
\NormalTok{engine }\OperatorTok{=}\NormalTok{ SCXDomainEngine(}\StringTok{"my\_domain"}\NormalTok{)}

\CommentTok{\# 3. 加载数据}
\NormalTok{X, y }\OperatorTok{=}\NormalTok{ load\_my\_domain\_data()}

\CommentTok{\# 4. 运行 SCX pipeline}
\NormalTok{engine.fit(X, y)}

\CommentTok{\# 5. 生成审计报告}
\NormalTok{report }\OperatorTok{=}\NormalTok{ engine.audit()}
\BuiltInTok{print}\NormalTok{(report.summary())}

\CommentTok{\# 6. 查看结果}
\BuiltInTok{print}\NormalTok{(report.n\_redundant, }\StringTok{"redundant samples"}\NormalTok{)}
\BuiltInTok{print}\NormalTok{(report.n\_noisy, }\StringTok{"noisy samples"}\NormalTok{)}
\BuiltInTok{print}\NormalTok{(report.n\_valuable, }\StringTok{"valuable samples"}\NormalTok{)}
\end{Highlighting}
\end{Shaded}

\subsubsection{12.3
新增领域清单}\label{ux65b0ux589eux9886ux57dfux6e05ux5355}

\begin{longtable}[]{@{}llll@{}}
\toprule\noalign{}
\# & 任务 & 责任人 & 工作量估计 \\
\midrule\noalign{}
\endhead
\bottomrule\noalign{}
\endlastfoot
1 & 实现 SCXStateEncoder 接口 & 领域开发者 & 0.5-2 天 \\
2 & 编写领域 YAML 配置 & 领域开发者 & 0.5 天 \\
3 & 准备测试数据 & 领域开发者 & 0.5-2 天 \\
4 & 运行测试 & 领域开发者 & 0.5 天 \\
5 & 注册到 DomainRegistry & 领域开发者 & 5 分钟 \\
\end{longtable}

总计：约 2-5 人天（根据领域复杂度）

\subsubsection{12.4 进阶：自定义 loss
function}\label{ux8fdbux9636ux81eaux5b9aux4e49-loss-function}

如果领域需要特殊损失函数（如力误差、物理约束违反），可以在 YAML 中指定：

\begin{Shaded}
\begin{Highlighting}[]
\CommentTok{\# domains/custom\_loss.yaml}
\FunctionTok{loss\_fn}\KeywordTok{:}
\AttributeTok{  }\FunctionTok{module}\KeywordTok{:}\AttributeTok{ my\_module.losses}
\AttributeTok{  }\FunctionTok{function}\KeywordTok{:}\AttributeTok{ my\_custom\_loss}
\AttributeTok{  }\FunctionTok{params}\KeywordTok{:}
\AttributeTok{    }\FunctionTok{alpha}\KeywordTok{:}\AttributeTok{ }\FloatTok{0.5}
\end{Highlighting}
\end{Shaded}

Adaapter 会自动加载并传入 ExpertReliability：

\begin{Shaded}
\begin{Highlighting}[]
\CommentTok{\# core/adapter.py 内部}
\ControlFlowTok{if} \StringTok{"loss\_fn"} \KeywordTok{in}\NormalTok{ domain\_config:}
\NormalTok{    loss\_cfg }\OperatorTok{=}\NormalTok{ domain\_config[}\StringTok{"loss\_fn"}\NormalTok{]}
\NormalTok{    module }\OperatorTok{=}\NormalTok{ importlib.import\_module(loss\_cfg[}\StringTok{"module"}\NormalTok{])}
\NormalTok{    loss\_fn }\OperatorTok{=} \BuiltInTok{getattr}\NormalTok{(module, loss\_cfg[}\StringTok{"function"}\NormalTok{])}
    \ControlFlowTok{if} \StringTok{"params"} \KeywordTok{in}\NormalTok{ loss\_cfg:}
\NormalTok{        loss\_fn }\OperatorTok{=}\NormalTok{ partial(loss\_fn, }\OperatorTok{**}\NormalTok{loss\_cfg[}\StringTok{"params"}\NormalTok{])}
\end{Highlighting}
\end{Shaded}

\begin{center}\rule{0.5\linewidth}{0.5pt}\end{center}

\subsection{13. 数据流全景}\label{ux6570ux636eux6d41ux5168ux666f}

\begin{verbatim}
用户输入
┌───────────────────────────────────────────────────────────────┐
│  X (原始数据)                                                  │
│  y (标签，可选)                                                │
│  experts (专家列表，可选)                                       │
└───────────────────┬───────────────────────────────────────────┘
                    │
                    ▼
┌───────────────────────────────────────────────────────────────┐
│  SCXDomainEngine.fit(X, y)                                    │
│                                                               │
│  1. Encoder.encode_batch(X) → Z (N, d)                       │
│                                                               │
│  2. Encoder.cluster(Z) → labels, centroids                    │
│                                                               │
│  3. Build StateSpace.from(labels, centroids)                  │
│     → ρ(s), D(s), C(s)                                        │
│                                                               │
│  4. ExpertRegistry.predict_all(X) → Y_pred (M, N)            │
│                                                               │
│  5. ExpertReliability.estimate(Y_pred, y, labels)            │
│     → R_m(s), SCX_m(s), uncertainties                        │
│                                                               │
│  6. DataClassifier.classify_all(state_metrics)               │
│     → Valuable / Redundant / Noisy / ExpertDependent         │
│                                                               │
│  7. ActionPolicy.decide(classifications, V(s))               │
│     → keep / compress / route / high-fidelity / ...          │
│                                                               │
│  8. SCXAuditReport.generate(...)                              │
│     → 结构化审计报告                                           │
└───────────────────┬───────────────────────────────────────────┘
                    │
                    ▼
┌───────────────────────────────────────────────────────────────┐
│  SCXAuditReport                                               │
│                                                               │
│  ├ State Map: 每个状态的质心、半径、密度、样本数                  │
│  ├ Redundancy: 冗余状态列表、压缩比例建议                         │
│  ├ Noise Risk: 噪声状态列表、建议动作                            │
│  ├ High Value: 高价值状态列表、补数据/训练建议                     │
│  ├ Expert Reliability Matrix: (M, K) SCX 可靠性分数             │
│  ├ Routing: 每个状态推荐的最佳专家                                │
│  └ Recommendations: 下一轮数据/仿真/标注预算建议                  │
└───────────────────────────────────────────────────────────────┘
\end{verbatim}

\begin{center}\rule{0.5\linewidth}{0.5pt}\end{center}

\subsection{14. 附录}\label{ux9644ux5f55}

\subsubsection{A.
当前代码与新架构的映射}\label{a.-ux5f53ux524dux4ee3ux7801ux4e0eux65b0ux67b6ux6784ux7684ux6620ux5c04}

\begin{longtable}[]{@{}
  >{\raggedright\arraybackslash}p{(\linewidth - 4\tabcolsep) * \real{0.3226}}
  >{\raggedright\arraybackslash}p{(\linewidth - 4\tabcolsep) * \real{0.3548}}
  >{\raggedright\arraybackslash}p{(\linewidth - 4\tabcolsep) * \real{0.3226}}@{}}
\toprule\noalign{}
\begin{minipage}[b]{\linewidth}\raggedright
当前文件
\end{minipage} & \begin{minipage}[b]{\linewidth}\raggedright
新架构归属
\end{minipage} & \begin{minipage}[b]{\linewidth}\raggedright
迁移动作
\end{minipage} \\
\midrule\noalign{}
\endhead
\bottomrule\noalign{}
\endlastfoot
\texttt{core/framework.py} & \texttt{core/} +
\texttt{core/domain\_engine.py} & 重构：SCXFramework 保持，新增
SCXDomainEngine \\
\texttt{core/config.py} & \texttt{core/config.py} & 扩展：支持从 YAML
加载领域配置 \\
\texttt{core/online.py} & \texttt{core/online.py} &
不变，也可作为插件 \\
\texttt{state/space.py} & \texttt{state/space.py} & 不变 \\
\texttt{expert/reliability.py} & \texttt{expert/reliability.py} &
不变 \\
\texttt{valuation/classifier.py} & \texttt{valuation/classifier.py} &
不变 \\
\texttt{action/policy.py} & \texttt{action/policy.py} & 不变 \\
\texttt{scx-health/src/scx\_health/encoder.py} &
\texttt{encoders/health.py} & 迁移并实现 SCXStateEncoder \\
\texttt{experiments/cifar/utils.py} & \texttt{encoders/vision.py} &
CIFAR 部分迁移，复用统一接口 \\
新增 & \texttt{encoders/base.py} & 新建 \\
新增 & \texttt{encoders/mlip.py} & 新建 \\
新增 & \texttt{encoders/trajectory.py} & 新建 \\
新增 & \texttt{core/domain\_engine.py} & 新建 \\
新增 & \texttt{core/adapter.py} & 新建 \\
新增 & \texttt{core/audit.py} & 新建 \\
新增 & \texttt{domains/registry.py} & 新建 \\
新增 & \texttt{domains/*.yaml} & 新建（每个领域一个） \\
\end{longtable}

\subsubsection{B.
完整用户入口示例}\label{b.-ux5b8cux6574ux7528ux6237ux5165ux53e3ux793aux4f8b}

```python \# ======== 方式 1：声明式（推荐） ======== from
scx.core.domain\_engine import SCXDomainEngine

\section{一行加载领域配置}\label{ux4e00ux884cux52a0ux8f7dux9886ux57dfux914dux7f6e}

engine = SCXDomainEngine(``scx/domains/mlip.yaml'')

\section{加载数据}\label{ux52a0ux8f7dux6570ux636e}

X, y = load\_dft\_data() \# 用户自己的数据加载

\section{运行完整 pipeline}\label{ux8fd0ux884cux5b8cux6574-pipeline}

engine.fit(X, y)

\section{获取结果}\label{ux83b7ux53d6ux7ed3ux679c}

report = engine.audit() print(report)

\section{压缩数据}\label{ux538bux7f29ux6570ux636e}

X\_sub, y\_sub = engine.compress(X, y, ratio=0.5)

\section{保存/加载}\label{ux4fddux5b58ux52a0ux8f7d}

engine.save(``outputs/my\_project.scx'') engine =
SCXDomainEngine.load(``outputs/my\_project.scx'')

\section{======== 方式 2：编程式（灵活）
========}\label{ux65b9ux5f0f-2ux7f16ux7a0bux5f0fux7075ux6d3b}

from scx.encoders.mlip import ACEEncoder from scx.core.framework import
SCXFramework from scx.core.config import SCXConfig

encoder = ACEEncoder(descriptor=``soap'', rcut=5.0) config =
SCXConfig(n\_states=20, n\_experts=3) framework =
SCXFramework(config=config) framework.fit(X, y, experts=expert\_list,
phi=encoder.encode) print(framework.summary())

\section{======== 方式 3：在线流式
========}\label{ux65b9ux5f0f-3ux5728ux7ebfux6d41ux5f0f}

from scx.core.online import OnlineSCXFramework import numpy as np

\section{初始质心}\label{ux521dux59cbux8d28ux5fc3}

centroids = np.random.randn(10, 128) online =
OnlineSCXFramework(centroids, M=3)

for x, expert\_id, loss in data\_stream(): record =
online.process\_sample(x, expert\_id, loss) if
record{[}``classification''{]} == ``noisy'': print(f''Alert: noisy
sample at state \{record{[}`state'{]}\}``)

\end{document}
